\documentclass[10pt]{article}

\usepackage[margin=0.85in]{geometry}
\usepackage{fontspec}
\setmainfont{TeX Gyre Termes}
\setsansfont{TeX Gyre Heros}
\setmonofont{Latin Modern Mono}
\usepackage{microtype}
\usepackage{xcolor}
\usepackage[table]{xcolor}
\usepackage{booktabs}
\usepackage{threeparttable}
\usepackage{tabularx}
\usepackage{array}
\usepackage[hidelinks]{hyperref}
\usepackage{url}
\usepackage{fvextra}

\newcommand{\best}[1]{\cellcolor{gray!15}\textbf{#1}}
\newcounter{algorithm}
\renewcommand{\thealgorithm}{S\arabic{algorithm}}
\renewcommand{\thesection}{S\arabic{section}}
\renewcommand{\thesubsection}{S\arabic{section}.\arabic{subsection}}
\renewcommand{\thetable}{S\arabic{table}}
\makeatletter
\renewcommand*\l@section{\@dottedtocline{1}{1.5em}{3em}}
\makeatother

\begin{document}

\begin{center}
{\Large Supplemental Material for}\\[3pt]
{\large A Controlled Study of HyDE-Style Query Expansion for Multi-Hop Evidence Acquisition}
\end{center}

\noindent
This supplement is organized as an auditable companion to the main manuscript. Tables~\ref{tab:configuration_provenance}--\ref{tab:efficiency_profile_full} provide the core lookup tables cited from the main text: configuration provenance, dataset artifact provenance, HotpotQA and 2Wiki paired statistics, deduplication sensitivity, guarded canonicalization decomposition, and full call/prompt-character accounting. Table~\ref{tab:baseline_implementation_parameters} lists baseline implementation details; Tables~\ref{tab:model_invocation_ledger}--\ref{tab:qwen_api_reproducibility} report model/API provenance; and Table~\ref{tab:corpus_scale_stress_full} gives the full corpus-scale retrieval stress test. Tables~\ref{tab:answer_metric_definition}--\ref{tab:supporting_evidence_metric_definition} define answer and evidence metric normalization before the answer-overlap diagnostics. Tables~\ref{tab:answer_overlap_definition}--\ref{tab:2wiki_hyde_query_mechanism} collect answer-overlap, gold-content, and 2Wiki query-composition diagnostics. Table~\ref{tab:query_length_matched_hyde} reports query-length matched sensitivity, and Table~\ref{tab:single_annotator_error_taxonomy} gives the single-annotator error taxonomy. Tables~\ref{tab:corpus_statistics}--\ref{tab:supporting_paragraph_audit} cover corpus statistics, supporting-title completeness, artifact availability, qualitative cases, and the supporting-paragraph audit. Table and section counters are independent: references always specify either \emph{Table Sx} or \emph{Section Sx}. Unless otherwise stated, artifact paths are relative to the anonymized project root.

\tableofcontents
\clearpage

\section{Configuration Provenance and Freezing Protocol}
\label{sec:supp_configuration_provenance}

The manuscript uses released IRCoT processed HotpotQA and 2WikiMultihopQA \texttt{test\_subsampled} splits under a processed local-corpus protocol. The following records separate retrieval-reader protocol choices from verifier-specific diagnostic choices by listing the selection source, the data inspected, and the artifact-level freeze evidence for each configuration. They do not claim that every design choice was selected using an independent development split. In particular, the $q+h$ row supports a narrower provenance claim: the final HyDE query composition was fixed before the reported $q$-only, $h$-only, and question-plus-rewrite diagnostic controls were generated and scored. Same-split diagnostic tuning is disclosed for the secondary verifier risk rules and guarded-merge policy.

\begin{table*}[t]
\centering
\caption{Configuration provenance for the reported processed-data evaluation.}
\label{tab:configuration_provenance}
\tiny
\begin{threeparttable}
\setlength{\tabcolsep}{2pt}
\renewcommand{\arraystretch}{0.92}
\begin{tabularx}{\textwidth}{>{\raggedright\arraybackslash}p{0.17\textwidth}
                            >{\raggedright\arraybackslash}p{0.27\textwidth}
                            >{\raggedright\arraybackslash}p{0.26\textwidth}
                            >{\raggedright\arraybackslash}X}
\toprule
Component & Selection basis & Data inspected & Frozen stage \\
\midrule
$q+h$ HyDE query composition & Chosen from the HyDE-style design rationale: keep the original question as an anchor and append a document-like hypothetical passage for dense retrieval. It is not claimed to be selected by optimizing over the reported $q$-only, $h$-only, or $q+\mathrm{rewrite}$ diagnostic rows. & Released processed-split schema and local retrieval artifacts were inspected to implement the pipeline. The reported $q$-only, $h$-only, and $q+\mathrm{rewrite}$ diagnostic results were generated after the final $q+h$ protocol was already in use. & Canonical $q+h$ HotpotQA retrieval and reader artifacts were frozen on 2026-07-04. Query-composition controls were added later: $q$-only and $h$-only on 2026-07-06, and $q+\mathrm{rewrite}$ on 2026-07-08. \\
HyDE prompt & Template asks for one concise hypothetical supporting passage for retrieval-query expansion. The design follows document-like HyDE-style expansion rather than a prompt search over final EM/F1. & Prompt-template design, released question text, and later generated-passage artifacts for answer-string-overlap auditing. Gold answers are retained for evaluation metadata but are not inserted into the prompt text. & Frozen before canonical HotpotQA HyDE generation. Prompt-builder and prompt-file hashes are reported in Table~\ref{tab:qwen_api_reproducibility}. \\
Rewrite prompt & Control prompt designed to spend one extra query-side Qwen3.7-Max call while producing a direct retrieval query rather than a document-like passage. & Same reported question set and control-design constraints. Gold answers are not inserted into prompt text. The prompt is a matched-call control, not an optimized alternative final method. & Frozen before query-reformulation generation and reader scoring on 2026-07-08 to 2026-07-09; prompt-builder hash is reported in Table~\ref{tab:qwen_api_reproducibility}. \\
Equal-budget query-form prompts & Post-hoc mechanism diagnostic added after the primary HotpotQA analysis to test whether document-like form remains useful at similar generated length and call budget. The four modes use a shared one-call, 40-word target, 35--45 word instruction window, and 64-token generation cap. & Primary HotpotQA results and the length-versus-form threat were inspected before designing this diagnostic. After the prompt family was fixed, the same templates and budgets were applied to HotpotQA and 2WikiMultihopQA without dataset-specific modification. Gold answers are retained as scoring metadata but are not inserted into prompt text. & Prompt JSONL files were generated on 2026-07-21. Full equal-budget Qwen3.7-Max generations were frozen on 2026-07-22 to 2026-07-23, followed by local BGE-base retrieval scoring on 2026-07-23. This row is a post-hoc robustness diagnostic, not part of the primary frozen study. \\
Reader prompt & Generic evidence-grounded short-answer prompt with common answer-type instructions for yes/no, location, date/age comparison, role/title, and answer-only formatting. & General multi-hop QA answer-format requirements and retrieved evidence serialization. The same reader prompt is used across retrieval variants and reused on 2WikiMultihopQA without dataset-specific retuning. The reader is not constrained to copy a contiguous evidence span. & Frozen before final reader prompt JSONL and answer artifacts were generated; prompt-builder hash is reported in Table~\ref{tab:qwen_api_reproducibility}. \\
Top-$k=10$ retrieval budget & Common evidence budget inherited from the processed-data retrieval-reader protocol and used for comparable retrieval metrics and reader prompts. & Released IRCoT processed-data format, local candidate-pool size, and prompt-length feasibility. & Applied uniformly to all reported methods and both datasets before final scoring; no method-specific top-$k$ tuning is reported. \\
900-character evidence truncation & Prompt-budget rule for comparable reader and verifier evidence blocks under hosted-model context limits. & Local corpus paragraph lengths and prompt serialization constraints were inspected during prompt construction. & Applied before final reader and verifier prompt JSONL artifacts were generated; the same limit is used for reader and verifier evidence blocks. \\
Verifier risk rules & Secondary verifier selection policy for likely format, alias, date, count, granularity, abstention, role/type, comparison, overlong-answer, and short-entity risks. & Same HotpotQA split diagnostics over canonical HyDE reader outputs and EM mismatches. This is disclosed as same-split verifier development, not as independent held-out tuning. & Frozen before final HotpotQA HyDE verifier prompt generation; applying the frozen policy selected 120 verifier cases. \\
Verifier prompt & Conservative JSON verifier prompt with \emph{keep} as default and a limited set of benchmark-style corrections. & Same HotpotQA verifier-development diagnostics and retrieved-evidence examples used to define safe correction categories. & Frozen before the 120 verifier calls; all 120 structured outputs are audited in Table~\ref{tab:qwen_api_reproducibility}. \\
Abstention guard & Deterministic merge rule to block a new abstention when the initial reader answer is not an abstention. & Same HotpotQA verifier diagnostics and prior harmful-abstention risk analysis. & Frozen before the final guarded merge. It is inactive on the final HyDE verifier outputs but retained as a preventative guard. \\
Numeric-unit guard & Deterministic merge rule to block unsupported additions of units to bare numeric answers. & Same HotpotQA verifier diagnostics and raw-verifier transition analysis, including observed correct-to-wrong numeric-unit expansions. & Frozen before the final guarded merge; the guard removes the observed correct-to-wrong transitions in the reported HyDE verifier reconstruction. \\
BGE-base stronger-encoder reproduction & Post-hoc robustness check asking whether the query-composition pattern persists with a stronger dense encoder. It is not used as a tuned replacement for the MiniLM primary protocol. & Primary MiniLM results and frozen generated query artifacts were inspected. The diagnostic changes the encoder to \texttt{BAAI/bge-base-en-v1.5} while keeping prompts, top-$k$, reader prompt, and local-corpus protocol fixed. & Added after the primary HotpotQA scoring pass; reported as a robustness reproduction rather than a pre-specified primary configuration. \\
Gold-aware scrubbing and answer-overlap diagnostics & Post-hoc internal-validity diagnostics for query-side answer-string overlap and generated gold-evidence traces. & Generated hypothetical passages, gold answers, supporting titles, and frozen retrieval artifacts were inspected for diagnostic grouping or scrubbing. These analyses are not used to choose a new deployed prompt. & Added after the primary HotpotQA results to bound the answer-leakage explanation; reported as diagnostic evidence only. \\
Corpus-scale and hard-negative audits & Post-hoc retrieval-only stress tests for corpus size, added-gold-title leakage, near-duplicate overlap, and harder processed-train distractors. & Released evaluation corpus records, processed-train distractors, frozen query-generation artifacts, and BGE similarity scores were inspected. The hard-negative row samples from the audited processed-train 100k pool rather than full Wikipedia. & Corpus-scale stress artifacts were generated after the primary study; leakage-filtered and BGE-hard audit artifacts are retained under the corpus-scale audit directory. These rows test external-validity threats after the primary study. \\
Qwen-Turbo fixed-evidence reader check & Post-hoc targeted reader robustness check. It changes only the reader model over fixed equal-budget top-10 evidence rather than rerunning query generation or retrieval. & Equal-budget retrieval outputs for direct rewrite, keyword/entity expansion, and document-like passage were inspected; no new retrieval artifacts were selected based on Qwen-Turbo answer quality. & Qwen-Turbo answer artifacts were frozen on 2026-07-23 to 2026-07-25. The check is a limited second-reader readout, not a broad reader sweep or primary end-to-end claim. \\
\bottomrule
\end{tabularx}
\begin{tablenotes}
\footnotesize
\item This table separates selection basis, inspected data, and freeze stage. It does not claim that every configuration was selected on an independent development split. The primary MiniLM retrieval-reader comparison is artifact-frozen as the main study; BGE, equal-budget, scrubbing, corpus-scale, top-$k$, supporting-paragraph, and Qwen-Turbo rows are post-hoc robustness diagnostics; same-split tuning is explicitly limited to the secondary verifier selection and guarded-merge analysis.
\end{tablenotes}
\end{threeparttable}
\end{table*}

\begin{table*}[t]
\centering
\caption{Dataset source and converted-artifact provenance for the processed local-corpus protocol.}
\label{tab:dataset_artifact_provenance}
\scriptsize
\begin{threeparttable}
\setlength{\tabcolsep}{3pt}
\newcommand{\shafull}[4]{\begin{tabular}{@{}l@{}}\texttt{#1}\\\texttt{#2}\\\texttt{#3}\\\texttt{#4}\end{tabular}}
\newcommand{\repocommit}{\begin{tabular}{@{}l@{}}\texttt{github.com/}\\\texttt{StonyBrookNLP/ircot}\\commit:\\\texttt{3c1820f698eea5eed}\\\texttt{db4fba3c56b64c96}\\\texttt{1e063e4}\\\texttt{main}; 2024-06-11\end{tabular}}
\begin{tabularx}{\textwidth}{>{\raggedright\arraybackslash}p{0.10\textwidth}
                            >{\raggedright\arraybackslash}p{0.22\textwidth}
                            >{\raggedright\arraybackslash}p{0.17\textwidth}
                            >{\raggedright\arraybackslash}p{0.20\textwidth}
                            >{\raggedright\arraybackslash}X}
\toprule
Dataset & Upstream repository and commit & Source file and local download & Raw SHA-256 & Converted JSONL SHA-256 \\
\midrule
HotpotQA &
\repocommit &
\begin{tabular}{@{}l@{}}\texttt{processed\_data/}\\\texttt{hotpotqa/}\\\texttt{test\_subsampled.jsonl}\\downloaded: 2026-06-30\\21:01:24\end{tabular} &
\shafull{729D8987560092C8}{476523488069E6FB}{9BD1527157DC0D7B}{FEBF3DDC336C558E} &
questions:
\shafull{23F4AD6F82182C1B}{FE34C7F42DA3E852}{125EA71B0B9138CC}{B3E1740B3501E8AA}
corpus:
\shafull{8B27434696F9AB47}{9EB7301B7CB51259}{FDFBD9C7CEA96D7A}{2A728B68F761ACB2} \\
\addlinespace
2Wiki &
\repocommit &
\begin{tabular}{@{}l@{}}\texttt{processed\_data/}\\\texttt{2wikimultihopqa/}\\\texttt{test\_subsampled.jsonl}\\downloaded: 2026-06-30\\21:01:21\end{tabular} &
\shafull{4948757B90627F40}{B77CD6BF81A62C1A}{7AD525D6A6A0C91B}{17A2E0E2FA88F5BC} &
questions:
\shafull{7E1871FE8813606A}{1E4FBE393A8A1C48}{CC67376CBB2E864F}{4F8926A37A3FF1C7}
corpus:
\shafull{CF7603B72D3B0BA0}{7F008C6FEEB281E}{2C63579098CF98F2}{C1E47E61F241069C0} \\
\bottomrule
\end{tabularx}
\begin{tablenotes}
\footnotesize
\item Converted artifacts are \texttt{questions.jsonl} and \texttt{corpus.jsonl} under \texttt{local-artifacts/datasets/ircot\_hotpotqa\_test500} and \texttt{local-artifacts/datasets/ircot\_2wikimultihopqa\_test500}. HotpotQA converted files have local timestamp 2026-06-30 21:17:21; 2WikiMultihopQA converted files have local timestamp 2026-07-04 22:15:50.
\item The conversion script is \texttt{src/data/prepare\_ircot\_hotpotqa.py}, SHA-256 prefix \texttt{C23D8CC9E6EB}. It reads each processed row's \texttt{contexts}, writes one question record with the first extracted answer and supporting-title list, and writes one paragraph-level corpus record per context with \texttt{doc\_id}, title, paragraph text, source question id, and \texttt{is\_supporting}. The local corpus is then formed as the within-dataset union of these question-level candidate pools.
\end{tablenotes}
\end{threeparttable}
\end{table*}


\clearpage
\section{Full HotpotQA Paired Statistics}
\label{sec:supp_hotpotqa_paired_statistics}

The HotpotQA paired statistics combine answer-side paired bootstrap/McNemar checks with retrieval-side paired checks over the same 500 examples. This table is the main supplemental lookup for the HotpotQA reliability claims in the manuscript.

\begin{table*}[t]
\centering
\caption{Full HotpotQA paired statistics on the released IRCoT \texttt{test\_subsampled} split.}
\label{tab:hotpotqa_paired_statistics}
\scriptsize
\begin{threeparttable}
\setlength{\tabcolsep}{3pt}
\begin{tabularx}{\textwidth}{>{\raggedright\arraybackslash}p{0.15\textwidth}
                            >{\raggedright\arraybackslash}X
                            r r c c}
\toprule
Endpoint & Baseline $\rightarrow$ target & $\Delta$ metric 1 [95\% CI] & $\Delta$ metric 2 [95\% CI] & W$\rightarrow$C / C$\rightarrow$W & Test $p$ \\
\midrule
Answer EM/F1 & One-step Dense RAG $\rightarrow$ HyDE-style RAG & 0.0900 [0.0620, 0.1200] & 0.0906 [0.0634, 0.1192] & 53 / 8 & $<0.0001$ \\
Answer EM/F1 & Single-query Reformulation RAG $\rightarrow$ HyDE-style RAG & 0.0880 [0.0600, 0.1180] & 0.0976 [0.0690, 0.1266] & 52 / 8 & $<0.0001$ \\
Answer EM/F1 & Question + rewritten query $\rightarrow$ Question + hypothetical passage & 0.1040 [0.0740, 0.1340] & 0.1057 [0.0791, 0.1364] & 59 / 7 & $2.38{\times}10^{-11}$ \\
Answer EM/F1 & BM25 + Dense Hybrid $\rightarrow$ HyDE-style RAG & 0.0440 [0.0160, 0.0740] & 0.0529 [0.0264, 0.0803] & 37 / 15 & 0.0032 \\
Answer EM/F1 & Rule-based Evidence-Expanded Iterative Retrieval $\rightarrow$ HyDE-style RAG & 0.0480 [0.0220, 0.0740] & 0.0540 [0.0296, 0.0788] & 35 / 11 & 0.0005 \\
Answer EM/F1 & HyDE-style RAG $\rightarrow$ HyDE-style RAG + Conservative Verifier & 0.0040 [0.0000, 0.0100] & 0.0030 [0.0000, 0.0080] & 2 / 0 & 0.5000 \\
\midrule
Retrieval All/Recall & Dense RAG $\rightarrow$ HyDE-style RAG & 0.2380 [0.2000, 0.2780] & 0.1200 [0.1010, 0.1410] & 120 / 1 & $<0.0001$ \\
Retrieval All/Recall & Single-query Reformulation RAG $\rightarrow$ HyDE-style RAG & 0.2500 [0.2140, 0.2860] & 0.1290 [0.1090, 0.1500] & 125 / 0 & $<0.0001$ \\
Retrieval All/Recall & Question + rewritten query $\rightarrow$ Question + hypothetical passage & 0.2400 [0.2020, 0.2780] & 0.1250 [0.1060, 0.1470] & 123 / 3 & $<0.0001$ \\
Retrieval All/Recall & Hybrid RAG $\rightarrow$ HyDE-style RAG & 0.1200 [0.0840, 0.1560] & 0.0590 [0.0410, 0.0780] & 76 / 16 & $<0.0001$ \\
Retrieval All/Recall & Rule-based Evidence-Expanded Iterative Retrieval $\rightarrow$ HyDE-style RAG & 0.1140 [0.0780, 0.1480] & 0.0580 [0.0390, 0.0780] & 72 / 15 & $<0.0001$ \\
\bottomrule
\end{tabularx}
\begin{tablenotes}
\footnotesize
\item All rows use the same 500 paired HotpotQA examples. For answer rows, metric 1 is $\Delta$EM, metric 2 is $\Delta$F1, W$\rightarrow$C/C$\rightarrow$W count exact-match answer transitions, and $p$ is the exact McNemar test over discordant EM outcomes. For retrieval rows, metric 1 is $\Delta$all-support hit@10, metric 2 is $\Delta$supporting-title recall@10, W$\rightarrow$C/C$\rightarrow$W count target-only and baseline-only full-support retrieval, and $p$ is the exact McNemar test over full-support retrieval transitions. Intervals are paired bootstrap 95\% percentile intervals over per-example deltas.
\end{tablenotes}
\end{threeparttable}
\end{table*}


\clearpage
\section{2Wiki Paired Statistics}
\label{sec:supp_2wiki_paired_statistics}

The 2WikiMultihopQA split is reported as a secondary generalization check under the same processed local-corpus protocol, reader prompt, top-$k=10$ retrieval budget, and local candidate-pool setting. The paired rows reuse frozen answer and retrieval artifacts and are interpreted as secondary diagnostics.

\begin{table*}[t]
\centering
\caption{2WikiMultihopQA paired statistics and secondary absolute results.}
\label{tab:2wiki_paired_statistics}
\scriptsize
\begin{threeparttable}
\setlength{\tabcolsep}{3pt}
\textit{Panel A: absolute retrieval-reader results.}
\begin{tabularx}{\textwidth}{>{\raggedright\arraybackslash}X r r r r r}
\toprule
Method & Any hit@10 & All hit@10 & Recall@10 & EM & F1 \\
\midrule
One-step Dense RAG & 0.9840 & 0.3580 & 0.6695 & 0.5400 & 0.5767 \\
BM25 RAG & \best{0.9980} & 0.4320 & 0.7250 & 0.5840 & 0.6169 \\
BM25 + Dense Hybrid & \best{0.9980} & 0.4480 & 0.7330 & 0.5840 & 0.6220 \\
HyDE-style RAG & 0.9940 & \best{0.6680} & \best{0.8575} & \best{0.6620} & \best{0.7289} \\
\bottomrule
\end{tabularx}

\vspace{5pt}
\textit{Panel B: paired HyDE-over-baseline diagnostics.}
\begin{tabularx}{\textwidth}{>{\raggedright\arraybackslash}X
                            r r c r r c}
\toprule
Baseline $\rightarrow$ target & $\Delta$EM [95\% CI] & $\Delta$F1 [95\% CI] & EM W$\rightarrow$C / C$\rightarrow$W & $\Delta$All@10 [95\% CI] & $\Delta$Recall@10 [95\% CI] & All@10 W$\rightarrow$C / C$\rightarrow$W \\
\midrule
Dense RAG $\rightarrow$ HyDE-style RAG & 0.1220 [0.0900, 0.1580] & 0.1523 [0.1186, 0.1881] & 72 / 11 & 0.3100 [0.2680, 0.3520] & 0.1880 [0.1665, 0.2085] & 157 / 2 \\
BM25 RAG $\rightarrow$ HyDE-style RAG & 0.0780 [0.0480, 0.1120] & 0.1121 [0.0793, 0.1450] & 54 / 15 & 0.2360 [0.1940, 0.2780] & 0.1325 [0.1115, 0.1550] & 133 / 15 \\
BM25 + Dense Hybrid $\rightarrow$ HyDE-style RAG & 0.0780 [0.0480, 0.1080] & 0.1069 [0.0751, 0.1399] & 53 / 14 & 0.2200 [0.1800, 0.2620] & 0.1245 [0.1020, 0.1475] & 125 / 15 \\
\bottomrule
\end{tabularx}
\begin{tablenotes}
\footnotesize
\item The 2WikiMultihopQA split is used as a secondary generalization check rather than as the main benchmark. Panel B uses the same 500 paired examples and frozen reader/retrieval artifacts. Answer-side intervals are paired bootstrap intervals over EM and F1 deltas; retrieval-side intervals are paired bootstrap intervals over all-support hit@10 and supporting-title recall@10 deltas. W$\rightarrow$C/C$\rightarrow$W count exact-match answer transitions for answer rows and target-only/baseline-only full-support retrieval for retrieval rows.
\end{tablenotes}
\end{threeparttable}
\end{table*}


\clearpage
\section{Deduplication Sensitivity}
\label{sec:supp_deduplication_sensitivity}

The local corpus for each dataset is the union of the released question-level candidate pools, indexed jointly within the dataset. Retrieval representations encode document titles and paragraph texts only. Question identifiers, supporting-evidence flags, gold answers, and evaluation metadata are retained for analysis but are not encoded by dense, BM25, or hybrid retrieval. Exact duplicate title--paragraph records are retained before indexing; if a paragraph appears in multiple question-level pools, each occurrence remains a separate document with a source-question-specific document identifier. Supporting-title metrics collapse duplicate retrieved titles before computing any-hit@10, all-hit@10, and supporting-title recall@10.

\begin{table*}[t]
\centering
\caption{Retrieval-side sensitivity to exact title--paragraph deduplication in the local joint index.}
\label{tab:dedup_retrieval_sensitivity}
\scriptsize
\begin{threeparttable}
\setlength{\tabcolsep}{4pt}
\begin{tabularx}{\textwidth}{>{\raggedright\arraybackslash}X l r r r r r r}
\toprule
Dataset & Method & Retain all@10 & Dedup all@10 & $\Delta$ all & Retain recall & Dedup recall & $\Delta$ recall \\
\midrule
HotpotQA & Dense & 0.6300 & 0.6320 & +0.0020 & 0.8080 & 0.8090 & +0.0010 \\
HotpotQA & Hybrid & 0.7480 & 0.7520 & +0.0040 & 0.8690 & 0.8710 & +0.0020 \\
HotpotQA & HyDE & 0.8680 & 0.8680 & +0.0000 & 0.9280 & 0.9280 & +0.0000 \\
2WikiMultihopQA & Dense & 0.3580 & 0.3920 & +0.0340 & 0.6695 & 0.6920 & +0.0225 \\
2WikiMultihopQA & Hybrid & 0.4480 & 0.4600 & +0.0120 & 0.7330 & 0.7415 & +0.0085 \\
2WikiMultihopQA & HyDE & 0.6680 & 0.6760 & +0.0080 & 0.8575 & 0.8625 & +0.0050 \\
\bottomrule
\end{tabularx}
\begin{tablenotes}
\footnotesize
\item Retain duplicates uses the canonical joint local index with repeated title--paragraph records preserved. Dedup removes exact normalized title--paragraph duplicates before retrieval, then rebuilds dense, BM25, hybrid, and HyDE retrieval artifacts with the same top-$k=10$ setting. Metrics still collapse duplicate supporting titles during evaluation. Reader answers are not regenerated for the deduplicated indexes, so EM/F1 are not reported in this table.
\item The CSV artifact additionally reports average duplicate title slots and exact duplicate slots in the retrieved top-10 lists.
\end{tablenotes}
\end{threeparttable}
\end{table*}


The sensitivity table shows that exact title--paragraph deduplication leaves the method ordering unchanged. HotpotQA HyDE metrics remain unchanged, while all three tested methods improve slightly on 2WikiMultihopQA where repeated local-pool records occupy more top-10 slots. Reader answers are not regenerated for the deduplicated indexes, so this analysis should be read as a retrieval-side robustness check rather than an EM/F1 robustness check. We therefore keep the retained-duplicate joint index as the canonical protocol because it mirrors the released local-pool artifacts, and treat deduplication as a robustness check rather than a protocol change.

\clearpage
\section{Guarded Answer Canonicalization Decomposition}
\label{sec:supp_verifier_guard_decomposition}

The verifier outputs are analyzed as a HyDE-native guarded answer-canonicalization layer. The decomposition below reuses the same HyDE reader answers and the same 120 risk-selected verifier outputs, then reconstructs raw merging and guarded merging policies without new model calls.

\begin{table*}[t]
\centering
\caption{HyDE verifier guard ablation reconstructed from existing verifier outputs.}
\label{tab:hyde_verifier_guard_ablation}
\small
\begin{threeparttable}
\begin{tabularx}{\textwidth}{>{\raggedright\arraybackslash}X r r r r r r c}
\toprule
Variant & $N$ & Verified & Changed & EM & F1 & $\Delta$F1 & W$\rightarrow$C / C$\rightarrow$W \\
\midrule
HyDE reader only & 500 & 0 & 0 & 0.6840 & 0.8105 & 0.0000 & 0 / 0 \\
HyDE + verifier raw & 500 & 120 & 5 & 0.6840 & 0.8119 & 0.0013 & 2 / 2 \\
HyDE + no-abstention guard & 500 & 120 & 5 & 0.6840 & 0.8119 & 0.0013 & 2 / 2 \\
HyDE + numeric-unit guard & 500 & 120 & 3 & 0.6880 & 0.8135 & 0.0030 & 2 / 0 \\
HyDE + both guards & 500 & 120 & 3 & 0.6880 & 0.8135 & 0.0030 & 2 / 0 \\
\bottomrule
\end{tabularx}
\begin{tablenotes}
\footnotesize
\item All rows reuse the same HyDE reader answers and the same 120 risk-selected verifier outputs. The raw row applies verifier final answers directly. The final row applies both deterministic guards: block new abstentions and block unsupported numeric-unit additions. The ablation is intended to justify the verifier's safety boundary, not to present verification as the main performance driver.
\end{tablenotes}
\end{threeparttable}
\end{table*}


\subsection{Guarded Merge Rule}
\label{sec:supp_guarded_merge_rule}

\refstepcounter{algorithm}
\noindent\textbf{Algorithm~\thealgorithm: Guarded canonicalization inference rule.}
\label{alg:guarded_canonicalization_rule}

\begin{enumerate}
    \item The reader receives the original question $q$ and retrieved corpus evidence $E$ and returns an initial short answer $a_0$.
    \item The rule-based selector computes a risk score $s(q,a_0)$. If $s(q,a_0)=0$, set the final answer $a_f=a_0$ and stop.
    \item If $s(q,a_0)>0$, query the conservative verifier with $(q,E,a_0)$ and parse its JSON output for the proposed final answer $\hat{a}$. If JSON parsing fails or the required answer field is missing, set $a_f=a_0$ and stop.
    \item Apply the abstention guard: if $a_0$ is not an abstention but $\hat{a}$ is an abstention, set $a_f=a_0$ and stop.
    \item Apply the numeric-unit guard: if $a_0$ is a bare number and $\hat{a}$ only appends a high-risk unit suffix, set $a_f=a_0$ unless the multi-word unit is explicitly requested by the question.
    \item Otherwise set $a_f=\hat{a}$.
\end{enumerate}
The final answer is therefore either the original reader answer or a selected-and-guarded verifier output for risk cases. The guard-decomposition table above reports the effect of raw merging and each deterministic guard using the same frozen verifier outputs.

\clearpage
\section{Full Call and Prompt-Character Accounting}
\label{sec:supp_full_accounting}

The main manuscript reports a compact LLM-call profile. The full accounting below separates query-side LLM calls, retrieval query executions, reader calls, verifier calls, total LLM calls, prompt-character proxies, evidence-character averages, and whether any method uses task-specific training. These are artifact-derived prompt-character and call-count statistics, not provider-side token, latency, or billing logs.

\begin{table*}[t]
\centering
\caption{Full structural call and retrieval-query accounting on the released IRCoT HotpotQA \texttt{test\_subsampled} split.}
\label{tab:efficiency_profile_full}
\scriptsize
\begin{threeparttable}
\setlength{\tabcolsep}{3pt}
\begin{tabularx}{\textwidth}{>{\raggedright\arraybackslash}X r r r r r r r}
\toprule
Method & Q-LLM/q & Ret.-q/q & Reader/q & Verifier/q & Total LLM/q & Prompt chars/q & Evidence chars/q \\
\midrule
LLM-only / No Retrieval & 0.00 & 0.00 & 1.00 & 0.00 & 1.00 & 368 & 0 \\
One-step Dense RAG & 0.00 & 1.00 & 1.00 & 0.00 & 1.00 & 6010 & 5149 \\
Multi-query RAG & 0.00 & 3.00 & 1.00 & 0.00 & 1.00 & 6010 & 5149 \\
Single-query Reformulation RAG & 1.00 & 1.00 & 1.00 & 0.00 & 2.00 & 6552 & 5169 \\
BM25 RAG & 0.00 & 1.00 & 1.00 & 0.00 & 1.00 & 6156 & 5295 \\
BM25 + Dense Hybrid & 0.00 & 2.00 & 1.00 & 0.00 & 1.00 & 5996 & 5135 \\
Rule-based Evidence-Expanded Iterative Retrieval & 0.00 & 4.00 & 1.00 & 0.00 & 1.00 & 6065 & 5204 \\
HyDE-style RAG & 1.00 & 1.00 & 1.00 & 0.00 & 2.00 & 6508 & 5121 \\
HyDE-style RAG + Guarded Canonicalization & 1.00 & 1.00 & 1.00 & 0.24 & 2.24 & 8214 & 5121 \\
\bottomrule
\end{tabularx}
\begin{tablenotes}
\footnotesize
\item Q-LLM/q counts query-side generation calls for matched-call reformulation or HyDE-style hypothetical-passage generation. Ret.-q/q counts retrieval query executions per question: Hybrid executes one dense and one BM25 retrieval, Multi-query executes three dense query variants, and Rule-based Evidence-Expanded Iterative Retrieval executes one first-round dense query plus three evidence-expanded dense queries. Total LLM/q counts query-side generation, reader answering, and selective verifier-backed canonicalization calls; it does not include retrieval query executions. Prompt chars/q reports average prompt characters per question from actual JSONL prompt files. Evidence chars/q reports serialized reader evidence length. The current logs do not store provider-side token usage, wall-clock latency, or billing. No method uses task-specific fine-tuning or model training.
\end{tablenotes}
\end{threeparttable}
\end{table*}


\begin{table*}[t]
\centering
\caption{Implemented baseline and control retrieval parameters used in the frozen HotpotQA artifacts.}
\label{tab:baseline_implementation_parameters}
\scriptsize
\begin{threeparttable}
\setlength{\tabcolsep}{3pt}
\begin{tabularx}{\textwidth}{>{\raggedright\arraybackslash}p{0.16\textwidth}
                            >{\raggedright\arraybackslash}p{0.27\textwidth}
                            >{\raggedright\arraybackslash}p{0.28\textwidth}
                            >{\raggedright\arraybackslash}X}
\toprule
Method & Query or document construction & Ranking and fusion rule & Frozen parameters and artifact notes \\
\midrule
BM25 RAG & Query tokens are extracted from the original question by lowercasing and matching \texttt{[a-z0-9]+}. Each corpus record is represented as \texttt{title title text}, so the title receives a two-count lexical boost and the paragraph body appears once. & Standard BM25 lexical scoring with Robertson-style IDF, term-frequency saturation, and document-length normalization. No stopword removal, stemming, phrase matching, synonym expansion, or learned reranking is applied. & $k_1=1.5$, $b=0.75$, top-$k=10$. The retained-duplicate local index is used for the canonical table rows; the deduplicated variant is reported only in Table~\ref{tab:dedup_retrieval_sensitivity}. \\
BM25 + Dense Hybrid & The dense input list is the one-step \texttt{all-MiniLM-L6-v2} top-10 list; the lexical input list is the BM25 top-10 list above. Both lists contain title and paragraph text only, with analysis metadata excluded from retrieval representations. & Reciprocal-rank fusion over the two ranked lists: $\mathrm{score}(d)=I[d\in D]/(60+r_D(d))+I[d\in B]/(60+r_B(d))$. Missing-list contributions are zero; fused documents are sorted by this score. & Dense weight $=1.0$, BM25 weight $=1.0$, RRF constant $=60$, dense candidates $=10$, BM25 candidates $=10$, final top-$k=10$. No dense/BM25 raw-score calibration is used. \\
Multi-query RAG & Three deterministic query variants are built from the question: the original question, a capitalized-or-numeric entity-term query, and a non-stopword keyword query. A punctuation-normalized fallback is used only if fewer than three distinct variants are available. No LLM is used. & Each query is encoded with the same dense encoder and retrieves 10 candidates. Candidate lists are fused by weighted reciprocal-rank fusion, with query-index weights $0.5^j$ in the frozen HotpotQA artifact and RRF constant $60$; ties are broken by the best weighted dense similarity stored for the candidate. & Number of queries $=3$, per-query candidates $=10$, final top-$k=10$, \texttt{query\_decay}$=0.5$. The frozen artifact has the same top-10 document set as one-step Dense RAG for all 500 HotpotQA examples, but only 39/500 examples have the same order; reader context order can therefore change without adding new documents. \\
Rule-based Evidence-Expanded Iterative Retrieval & The first stage is the one-step dense top-10 retrieval. The second stage constructs four dense queries: the original question plus one expanded query for each of the first three first-round documents, formatted as the question, the related evidence title, and up to 700 characters of related evidence text. No reader output and no LLM call are used to form these queries. & The four dense queries each retrieve 10 candidates. Candidate documents are merged by document ID; each candidate keeps its maximum adjusted dense similarity, with weight $1.0$ for the original query and $0.97$ for evidence-expanded queries. The final list is sorted by the retained adjusted score. & First-round expansion documents $=3$, per-query candidates $=10$, retrieval query executions per question $=4$, final top-$k=10$. This is a fixed heuristic iterative control rather than a reproduction of IRCoT. The final top-10 is a re-ranked union of the second-stage candidate set, not a simple append to the first-stage list. \\
\bottomrule
\end{tabularx}
\begin{tablenotes}
\footnotesize
\item These rows describe the implemented retrieval baselines and controls used for the reported HotpotQA comparisons. Dense retrieval, Multi-query retrieval, Rule-based Evidence-Expanded Iterative Retrieval, and HyDE-style retrieval all use the same \texttt{sentence-transformers/all-MiniLM-L6-v2} encoder and normalized FAISS inner-product search unless a row states otherwise. Multi-query and Rule-based Evidence-Expanded Iterative Retrieval are fixed heuristic controls, and Single-query Reformulation is a matched-call mechanism control; they are not claimed as faithful public-method reproductions. The matched-call Single-query Reformulation and HyDE-style query prompts are specified in Section~\ref{sec:supp_prompt_api}.
\end{tablenotes}
\end{threeparttable}
\end{table*}


\clearpage
\section{Prompt and API Configuration}
\label{sec:supp_prompt_api}

This section gives the complete prompt templates used by the canonical query-side generation, matched-call query reformulation, equal-budget query-form diagnostic, evidence-grounded short-answer reader, and conservative verifier. The reader prompt uses generic short-answer formatting instructions for common multi-hop QA answer types and is applied unchanged to HotpotQA and 2WikiMultihopQA. The reader and verifier evidence blocks serialize the top-10 retrieved passages with title and paragraph text, truncated to at most 900 characters per document.

\subsection{HyDE Hypothetical-Passage Generation Prompt}
\begingroup\footnotesize
\begin{verbatim}
You are generating a hypothetical supporting passage for retrieval.
Given the question below, write a concise factual-sounding supporting passage
that would likely contain the answer.
The passage may be imperfect, but it should include key entities, relations,
and context useful for retrieving evidence.
Do not explain the task. Do not include bullet points. Write one short paragraph.

Question:
{question}

Hypothetical supporting passage:
\end{verbatim}
\endgroup

\subsection{Matched-Call Single-Query Reformulation Prompt}
\begingroup\footnotesize
\begin{verbatim}
You are rewriting a multi-hop question into a single retrieval query.
Given the question below, produce one concise search query for dense retrieval.
Use only entities, relations, and constraints already present in the question.
Do not answer the question. Do not add facts that are not stated in the question.
Do not write a paragraph or explanation. Output only the single retrieval query.

Question:
{question}

Single retrieval query:
\end{verbatim}
\endgroup

\subsection{Equal-Budget Query-Form Prompts and Serialization}
\label{sec:supp_equal_budget_prompts}

The equal-budget diagnostic reported in main Tables 5 and 6 uses four separate query-form prompts with a shared wrapper. Each prompt asks for retrieval text, not an answer, uses one query-side generation call, targets 35--45 English words around a 40-word target, and sets the generation budget to \texttt{max\_tokens=64} with temperature 0. The prompt rows store gold answers and supporting-fact metadata only for later scoring; these fields are not inserted into the user prompt. We did not discard, truncate, or regenerate outputs solely because they fell outside the requested word window; the realized word counts are reported with the retrieval summary. Failed API requests are retried by the generation wrapper and, under \texttt{--resume}, rerun only as missing rows rather than replaced by hand-written text.

The generated text is used for retrieval only. For all four equal-budget rows, the BGE-base encoder input is the original question plus the generated retrieval text:
\begingroup\footnotesize
\begin{verbatim}
{question}

Generated retrieval text:
{generated_text}
\end{verbatim}
\endgroup
Thus, the document-like condition is encoded as $[q;G_{\mathrm{doc}}(q)]$, not as $G_{\mathrm{doc}}(q)$ alone. The exact separator is two newline characters, the label ``Generated retrieval text:'', and one newline before the model output. The BGE-base tokenizer then encodes this single serialized string with a 512-token query limit and right truncation. Lists, semicolons, bullets, and sentences are controlled only by the prompt text below: the keyword mode requests comma-separated phrases rather than sentences, the direct-rewrite mode requests one sentence, the decomposition mode requests numbered subquestions, and the document-like mode requests prose rather than bullets. The prompts do not include retrieved evidence or corpus text. They also do not impose a global ban on adding likely bridge concepts: the keyword/entity and document-like variants explicitly invite likely bridge concepts or plausible supporting context, while the direct-rewrite variant asks only to preserve the original intent and make implicit relations explicit.

\subsubsection*{Keyword/Entity Expansion}
\begingroup\footnotesize
\begin{Verbatim}[breaklines=true,breakanywhere=false,breaksymbolleft={},breaksymbolright={}]
You are generating retrieval text for a controlled multi-hop QA query-composition diagnostic.
Write a compact keyword/entity list for dense retrieval. Include entities, relations, constraints, and likely bridge concepts. Use comma-separated phrases, not sentences.

Use only one generation call. Do not answer the question. Do not explain the task.
Aim for 35-45 English words. Prefer concrete entities and relations over filler.
Output only the retrieval text.

Question:
{question}

Retrieval text:
\end{Verbatim}
\endgroup

\subsubsection*{Direct Rewrite}
\begingroup\footnotesize
\begin{Verbatim}[breaklines=true,breakanywhere=false,breaksymbolleft={},breaksymbolright={}]
You are generating retrieval text for a controlled multi-hop QA query-composition diagnostic.
Write one single dense-retrieval rewrite of the question. Preserve the original intent, make implicit relations explicit, and keep it as one sentence.

Use only one generation call. Do not answer the question. Do not explain the task.
Aim for 35-45 English words. Prefer concrete entities and relations over filler.
Output only the retrieval text.

Question:
{question}

Retrieval text:
\end{Verbatim}
\endgroup

\subsubsection*{Question Decomposition}
\begingroup\footnotesize
\begin{Verbatim}[breaklines=true,breakanywhere=false,breaksymbolleft={},breaksymbolright={}]
You are generating retrieval text for a controlled multi-hop QA query-composition diagnostic.
Write numbered subquestions that decompose the multi-hop question into retrieval steps. Keep the total output within the same length budget.

Use only one generation call. Do not answer the question. Do not explain the task.
Aim for 35-45 English words. Prefer concrete entities and relations over filler.
Output only the retrieval text.

Question:
{question}

Retrieval text:
\end{Verbatim}
\endgroup

\subsubsection*{Document-Like Passage}
\begingroup\footnotesize
\begin{Verbatim}[breaklines=true,breakanywhere=false,breaksymbolleft={},breaksymbolright={}]
You are generating retrieval text for a controlled multi-hop QA query-composition diagnostic.
Write one short encyclopedia-style passage that could plausibly support answering the question. Use prose rather than bullets.

Use only one generation call. Do not answer the question. Do not explain the task.
Aim for 35-45 English words. Prefer concrete entities and relations over filler.
Output only the retrieval text.

Question:
{question}

Retrieval text:
\end{Verbatim}
\endgroup

\subsection{Evidence-Grounded Short-Answer Reader Prompt}
The frozen prompt text below uses the historical word ``extractive'' in its system role. In this manuscript, this denotes evidence-constrained short-answer reading rather than standard span-extractive QA: the answer is not required to be a contiguous span copied from one evidence passage.
\begingroup\footnotesize
\begin{verbatim}
You are an extractive multi-hop question answering system.
Use only the evidence below. Combine facts across evidence items when needed.
Return the exact short answer phrase requested by the question.

Important rules:
- Return only the answer, with no explanation.
- Prefer the most specific answer phrase in the evidence.
- For yes/no questions, answer only "yes" or "no".
- For location questions, include the full location if the evidence gives it.
- For date/age comparison questions, the earlier birth date means the person is older.
- For role/position questions, return the exact position title, not a broader category.
- Do not answer "I don't know" if the evidence contains enough information to infer the answer.

Question:
{question}

Evidence:
{evidence}

Exact short answer:
\end{verbatim}
\endgroup

\subsection{Conservative Verification Prompt}
\begingroup\footnotesize
\begin{verbatim}
You are a conservative evidence-grounded answer verification agent.

Your default behavior is to KEEP the initial answer. Rewrite it only when there
is a clear benchmark-style formatting issue and the evidence directly supports
the corrected phrase.

Allowed corrections:
1. Nationality questions: prefer adjectival nationality when the evidence supports it.
   Example: "American" instead of "United States citizen".
2. Count questions: if the initial answer is only a number but the question asks
   for a counted object, include the unit from the evidence. Example:
   "12 member universities" instead of "12".
3. Date questions: if the question asks for an event time and the evidence gives
   a more complete month/year phrase, include the complete phrase.
4. Yes/no questions: if the question asks a yes/no comparison and the evidence
   directly supports it, normalize the answer to "Yes" or "No".
5. Abstained answers: if the initial answer is "I don't know" but the evidence
   directly contains the exact requested short answer, provide that answer.
   Otherwise keep "I don't know".
6. Location, role, title, or type questions: only add a missing qualifier when
   the question asks for that qualifier and the exact phrase appears in evidence.

Do not make these changes:
1. Do not abstain if the initial answer is directly supported by any evidence.
2. Do not expand a familiar short person name to a full legal name unless the
   question explicitly asks for the full name.
3. Do not remove location qualifiers such as state/country names.
4. Do not remove parenthetical title explanations.
5. Do not rewrite a correct answer merely because another phrase is also present.
6. Do not add units to a numeric answer when the question only asks for a number.
7. Do not replace a precise answer with a longer phrase unless the longer phrase
   is explicitly requested by the question.

Return only valid JSON with this exact schema:
{"verdict":"keep|correct|abstain","final_answer":"...","reason":"short reason"}

Question:
{question}

Initial answer:
{prediction}

Evidence:
{evidence}

JSON:
\end{verbatim}
\endgroup

\subsection{API, Model Invocation, and Result Ledger}

The hosted model is recorded at the artifact level because the provider exposes a hosted alias rather than a public immutable checkpoint. Reported reader, query-generation, and verifier artifacts use the stored API model string \texttt{qwen3.7-max}. The following tables freeze the local invocation windows, stored model fields, prompt hashes, retry policy, parsing policy, and raw-output boundaries.

\begin{table*}[t]
\centering
\caption{Model invocation ledger for frozen hosted-Qwen artifacts.}
\label{tab:model_invocation_ledger}
\scriptsize
\begin{threeparttable}
\setlength{\tabcolsep}{3pt}
\begin{tabularx}{\textwidth}{>{\raggedright\arraybackslash}p{0.18\textwidth}
                            >{\raggedright\arraybackslash}p{0.24\textwidth}
                            >{\raggedright\arraybackslash}p{0.16\textwidth}
                            r
                            >{\raggedright\arraybackslash}X}
\toprule
Artifact group & Local freeze window & Stored model field(s) & Calls / rows & Frozen artifact source \\
\midrule
HotpotQA retrieval baselines & 2026-07-02 to 2026-07-04 & \texttt{qwen3.7-max} & 5 $\times$ 500 & LLM-only, Dense, BM25, Hybrid, and Multi-query answer JSONL files. \\
HotpotQA matched-call query reformulation & 2026-07-08 to 2026-07-09 & \texttt{qwen3.7-max} & 3 $\times$ 500 & Single-query reformulation generation, rewritten-query reader, and question + rewritten-query reader JSONL files. \\
HotpotQA HyDE and iterative runs & 2026-07-04 to 2026-07-06 & \texttt{qwen3.7-max} & 4 $\times$ 500 & HyDE generation, HyDE reader, Rule-based Evidence-Expanded Iterative Retrieval reader, and HyDE hypothetical-only mechanism answer files. \\
HotpotQA conservative verifier & 2026-07-04 & \texttt{qwen3.7-max} & 120 verifier calls; 500 merged rows & Risk-selected verifier answer JSONL plus numeric-guarded merged evaluation file. \\
2WikiMultihopQA secondary runs & 2026-07-05 to 2026-07-06 & \texttt{qwen3.7-max} & 5 $\times$ 500 & 2Wiki Dense, BM25, Hybrid, HyDE generation, and HyDE reader answer files. \\
Equal-budget query-form generation & 2026-07-21 to 2026-07-23 & \texttt{qwen3.7-max} & 8 $\times$ 500 & Post-hoc HotpotQA and 2Wiki keyword/entity, direct-rewrite, question-decomposition, and document-like generation files. \\
Equal-budget Qwen-Turbo reader check & 2026-07-23 to 2026-07-25 & \texttt{qwen-turbo} & 6 $\times$ 500 & Direct-rewrite, keyword/entity-expansion, and document-like top-10 reader answer files for HotpotQA and 2Wiki. \\
\bottomrule
\end{tabularx}
\begin{tablenotes}
\footnotesize
\item All rows use the Alibaba Cloud Bailian/DashScope OpenAI-compatible chat-completions API wrapper with temperature-0 decoding (\texttt{temperature=0.0}). The primary pipeline uses the stored API model string \texttt{qwen3.7-max}; the equal-budget query-form generation reuses \texttt{qwen3.7-max} as a post-hoc diagnostic; the targeted cross-reader check intentionally uses \texttt{qwen-turbo}. The configured \texttt{max\_tokens} budgets are reported separately as 64 for HyDE passage generation, equal-budget query-form generation, rewritten-query generation, reader answers, and verifier JSON output. Temperature 0 does not guarantee byte-identical replay across requests or hosted-model alias updates. Freeze windows are local artifact timestamps, not provider request logs. Table~\ref{tab:qwen_api_reproducibility} records the API-level reproducibility boundary, and the complete artifact-level ledger is recorded in \texttt{docs/model\_invocation\_ledger.md}.
\end{tablenotes}
\end{threeparttable}
\end{table*}

\begin{table*}[t]
\centering
\caption{Hosted Qwen API reproducibility record for frozen generation artifacts.}
\label{tab:qwen_api_reproducibility}
\scriptsize
\begin{threeparttable}
\setlength{\tabcolsep}{3pt}
\begin{tabularx}{\textwidth}{>{\raggedright\arraybackslash}p{0.20\textwidth}
                            >{\raggedright\arraybackslash}X}
\toprule
Field & Frozen record or limitation \\
\midrule
Provider and interface & Alibaba Cloud Bailian/DashScope OpenAI-compatible chat-completions interface, accessed through the local \texttt{LLM\_BASE\_URL} environment variable. API credentials and endpoint URL are not stored in the artifact release. \\
API model string & Reported primary reader, query-generation, and verifier artifacts use the hosted API model string \texttt{qwen3.7-max}. The post-hoc equal-budget query-form generation also uses \texttt{qwen3.7-max}, while the targeted equal-budget reader robustness check uses \texttt{qwen-turbo} as a second hosted reader and is reported separately. The model string is stored in the local JSONL artifacts when the generation wrapper writes model metadata. \\
Snapshot/version & No immutable provider-side snapshot identifier was available for the hosted \texttt{qwen3.7-max} alias. Canonical runs are therefore frozen by local prompt and output artifacts rather than by a public model checkpoint. \\
Module model-string consistency & Primary HyDE generation, matched-call query reformulation, reader, 2Wiki, and verifier runs use the same API model string, \texttt{qwen3.7-max}. The post-hoc equal-budget query-form generation also uses \texttt{qwen3.7-max} but is not part of the primary frozen scoring window. The Qwen-Turbo check intentionally changes only the reader over fixed equal-budget retrieval evidence. Because these are hosted aliases rather than immutable snapshots, the paper does not claim replay against the exact same provider-internal model state after an alias evolves. \\
Experiment window & Primary Qwen3.7-Max artifacts were generated or frozen between 2026-07-02 and 2026-07-09 local time. Post-hoc equal-budget Qwen3.7-Max query-form generations were frozen between 2026-07-21 and 2026-07-23 local time. The targeted Qwen-Turbo fixed-evidence reader artifacts were frozen between 2026-07-23 and 2026-07-25 local time. Artifact-level windows and paths are listed in Table~\ref{tab:model_invocation_ledger}. \\
Endpoint region/type & Endpoint type is OpenAI-compatible \texttt{/v1} chat completions. Provider region was not exported into the local JSONL artifacts and is therefore not claimed as a reproducible field. \\
Temperature and output budgets & \texttt{temperature=0.0}. Configured \texttt{max\_tokens}: HyDE passage generation = 64; rewritten-query generation = 64; reader answer generation = 64; verifier JSON generation = 64. Temperature 0 is used to reduce sampling variation, but byte-identical replay is not claimed for a hosted alias. \\
\texttt{top\_p}, seed, and thinking mode & No \texttt{top\_p}, seed, logprob, response-format, or reasoning/thinking-mode parameter is explicitly sent by the wrapper, so provider defaults apply where relevant. Seed-controlled replay and thinking-mode equivalence are therefore not claimed. \\
Retry and timeout policy & Timeout is 60 seconds by default. Each request is attempted up to 3 times with 5 seconds between retries. If all retries fail, the script raises the final exception rather than writing a synthetic answer row. The \texttt{--resume} option skips completed IDs already present in the output JSONL and appends only missing rows. \\
API failure replacement policy & Reported \texttt{qwen3.7-max} HyDE, matched-call, 2Wiki, reader, and verifier runs, as well as the \texttt{qwen-turbo} fixed-evidence reader check, use completed rows from the frozen output files. Failed requests are rerun only as missing rows under \texttt{--resume}; they are not silently replaced by hand-written answers. \\
Prompt version/hash & Prompt-builder script SHA-256 prefixes: HyDE \texttt{89ECB4FB}, reader \texttt{98904C37}, query reformulation \texttt{FF9EA7EF}, equal-budget query-form \texttt{A234A00B}, verifier \texttt{BB2495BC}. Canonical HotpotQA prompt-file SHA-256 prefixes: HyDE generation \texttt{4D43047C}, HyDE reader \texttt{DCBC0B3E}, verifier risk-120 \texttt{763D188C}. Equal-budget prompt-file SHA-256 prefixes: HotpotQA keyword \texttt{048B13BE}, direct \texttt{2777A456}, decomposition \texttt{4244BD3E}, document-like \texttt{FF4A4ECF}; 2Wiki keyword \texttt{EC679B8B}, direct \texttt{E7464644}, decomposition \texttt{BF51AE1A}, document-like \texttt{CF15FFFD}. \\
Malformed JSON handling & Reader, HyDE, and query-reformulation generations are stored as raw text predictions and are not JSON-parsed. Verifier outputs are parsed as JSON when possible, with substring JSON extraction for wrapped responses. If verifier parsing fails, or if the parsed object lacks a nonempty \texttt{final\_answer} field, the fail-safe policy retains the initial reader answer $a_0$; malformed verifier text is never used as a candidate final answer. No automatic re-request is triggered solely by malformed JSON or JSON parse failure. \\
Verifier JSON parse audit & The canonical HotpotQA HyDE verifier artifact contains 120 verifier rows generated with \texttt{max\_tokens=64}. All 120 \texttt{prediction} fields were successfully parsed as direct JSON objects; substring fallback was used for 0 rows; final JSON parse failures were 0; empty verifier outputs were 0. All 120 verifier rows store the model field \texttt{qwen3.7-max}. The wrapper does not retain provider \texttt{finish\_reason} fields or per-request retry counters, and no automatic re-request rule is triggered solely by suspected output truncation. \\
Raw request/response retention & Local artifacts retain the full user prompt text, retrieved evidence, prediction text, model field, gold-answer metadata for scoring, and supporting-fact metadata. They do not retain the provider-side raw request envelope, raw response envelope, request IDs, token accounting, latency, billing, or per-request provider timestamps. Local artifact modification times provide freeze timestamps in the ledger. \\
\bottomrule
\end{tabularx}
\begin{tablenotes}
\footnotesize
\item Because the hosted model alias may evolve, all reported generations are accompanied by frozen raw prompts, local raw text outputs, local artifact freeze timestamps, and stored model strings. This table supports artifact-level reproducibility and rescoring, not exact replay of a provider-internal model snapshot or provider-side request log reconstruction.
\end{tablenotes}
\end{threeparttable}
\end{table*}


For the structured verifier output, we audited the canonical HotpotQA HyDE verifier artifact
listed in Table~\ref{tab:model_invocation_ledger}.
All 120 verifier outputs were successfully parsed as direct JSON under the reported 64-token
budget. No substring JSON-recovery fallback was needed, no empty verifier output was observed,
and no final JSON parse failure remained in the frozen verifier artifact. The generation wrapper
does not store provider-side \texttt{finish\_reason} or per-request retry counters, so truncation
is bounded here by artifact-level parse success rather than by provider response metadata.

\subsection{Dense Encoder and Truncation Configuration}
\label{sec:supp_dense_encoder_truncation}

All canonical MiniLM-based dense retrieval variants use \path{sentence-transformers/all-MiniLM-L6-v2} through the \path{SentenceTransformer.encode} interface with normalized embeddings and FAISS inner-product search. The local audit environment reports \path{sentence-transformers==5.6.0}; the local MiniLM model cache records a \path{sentence_bert_config.json} value of \path{max_seq_length=256}. The underlying uncased BERT tokenizer uses WordPiece tokenization and the default Transformers truncation side is right truncation. Thus, MiniLM dense inputs are encoded as a single sequence with a shared 256-WordPiece-token budget, including the tokenizer's special tokens, and excess content is dropped from the right after serialization. The retained HotpotQA and 2Wiki MiniLM indexes are FAISS \path{IndexFlatIP} indexes over normalized 384-dimensional embeddings, so inner-product search is cosine-style nearest-neighbour retrieval.

The same tokenizer, 256-token budget, and right-truncation rule are applied to all dense query-composition diagnostic rows on HotpotQA: question only ($q$), single rewritten query ($r$), question plus rewritten query ($q+r$), hypothetical passage only ($h$), and question plus hypothetical passage ($q+h$). The generated text is character-capped before this tokenizer step: reformulated queries use at most 300 characters, and HyDE hypothetical passages use at most 900 characters with an ellipsis appended if capped. The combined rows share one encoder input sequence rather than separate per-field budgets for the question and generated text.

Dense query serialization is fixed by the retrieval scripts. The $q+h$ row serializes the original question, two newline characters, the label ``Hypothetical supporting passage:'', one newline, and the generated passage. The $q+r$ row uses the same layout with the label ``Rewritten retrieval query:''. The $q$, $r$, and $h$ rows encode only their corresponding serialized text. Corpus documents are serialized for dense indexing as \path{{title}. {paragraph}} and use the same 256-token right-truncation rule; there is no separate title or paragraph token budget. Question identifiers, support flags, answers, and evaluation metadata are not encoded.

All post-hoc BGE-base diagnostics use one frozen encoder configuration across the rows being compared:
\begin{itemize}
\item model identifier: \path{BAAI/bge-base-en-v1.5};
\item local Hugging Face \path{main} revision: \path{a5beb1e3e68b9ab74eb54cfd186867f64f240e1a};
\item local \path{model.safetensors} SHA-256 prefix: \path{C7C1988AAE20};
\item local \path{config.json} SHA-256 prefix: \path{BC00AF31A4A3};
\item local \path{tokenizer_config.json} SHA-256 prefix: \path{9261E7D79B44};
\item implementation: \path{transformers.AutoTokenizer} and \path{AutoModel}, not \path{SentenceTransformer.encode}; the local audit environment reports \path{transformers==4.57.6} and \path{torch==2.8.0}.
\end{itemize}

The BGE diagnostics use CLS pooling, taking \path{last_hidden_state[:, 0]} as the sequence embedding. Corpus-scale scripts expose a mean-pooling option for sensitivity work, but the reported BGE rows use the default \path{pooling=cls}. Both document and query inputs use \path{max_length=512}, tokenizer-side right truncation, and identical L2 normalization with \path{F.normalize(..., p=2, dim=1)}; this is equivalent to using normalized embeddings. No query instruction prefix was added. Queries and documents are encoded as plain serialized strings, without BGE's recommended instruction-style query prefix.

BGE document serialization is \path{{title}. {paragraph_text}} with no separate title budget. The stronger-encoder query-composition diagnostic serializes question-only as $q$, single rewrite as $r$, hypothetical-only as $h$, and uses the following combined-query strings:
\begingroup\footnotesize
\begin{verbatim}
{question}\n\nRewritten retrieval query:\n{rewrite}

{question}\n\nHypothetical supporting passage:\n{hypothetical passage}
\end{verbatim}
\endgroup
The equal-budget diagnostic serializes each generated form identically as:
\begingroup\footnotesize
\begin{verbatim}
{question}\n\nGenerated retrieval text:\n{generated_text}
\end{verbatim}
\endgroup
Thus the document-like condition encodes $[q;G_{\mathrm{doc}}(q)]$ rather than $G_{\mathrm{doc}}(q)$ alone. BGE corpus-scale and hard-negative diagnostics use the same question-only, single-rewrite, and question-plus-hypothetical serializations as the stronger-encoder diagnostic. BGE diagnostics do not build a FAISS index: they compute dense score matrices as the dot product between L2-normalized query and document embeddings and select top-$k$ with \path{torch.topk}, making the score equivalent to cosine similarity.

\section{Corpus-Scale Retrieval Stress Test}
\label{sec:supp_corpus_scale_stress}

The main experiments use processed local candidate pools. To test sensitivity to a substantially noisier index without claiming full-wiki retrieval, we compare the original retained-duplicate local index with expanded 50k and 100k stress-test indexes. The expanded corpora preserve the released evaluation corpus and then add exact-title--paragraph-deduplicated distractor passages sampled from the corresponding IRCoT processed training split using seed 20260721. The expansion is gold-preserving by construction because the original evaluation support passages remain in the index; it is therefore a diagnostic retrieval stress test rather than a deployable retrieval protocol.

All rows use \texttt{BAAI/bge-base-en-v1.5}, CLS pooling, 512-token dense inputs, top-$k=10$, and the frozen question-only, single-query reformulation, and HyDE hypothetical-passage artifacts. No reader calls are rerun for this stress test. The script is \texttt{src/evaluation/analyze\_corpus\_scale\_retrieval\_stress.py}; the summary CSV, expanded corpora, retrieval JSONL files, and embedding caches are stored under \texttt{local-artifacts/corpus\_scale\_stress}. Paired bootstrap intervals are computed over per-example retrieval deltas.

\begin{table*}[t]
\centering
\caption{Full corpus-scale retrieval stress-test results. Local rows reuse the retained-duplicate BGE local-index artifacts from the main stronger-encoder reproduction table; only the 50k and 100k rows use newly expanded stress-test corpora.}
\label{tab:corpus_scale_stress_full}
\scriptsize
\begin{threeparttable}
\setlength{\tabcolsep}{2.5pt}
\begin{tabularx}{\textwidth}{l r >{\raggedright\arraybackslash}X r r r >{\raggedright\arraybackslash}X >{\raggedright\arraybackslash}X}
\toprule
Dataset & Docs & Query input & Any & All & Recall & HyDE delta vs. question & HyDE delta vs. rewrite \\
\midrule
HotpotQA & Local & Question only & 0.9940 & 0.8820 & 0.9380 & -- & -- \\
HotpotQA & Local & Single rewritten query & 0.9920 & 0.8340 & 0.9130 & -- & -- \\
HotpotQA & Local & Question + hypothetical passage & 0.9940 & \best{0.9200} & \best{0.9570} & All +0.0380 [0.0100, 0.0660]; recall +0.0190 [0.0050, 0.0340] & All +0.0860 [0.0560, 0.1160]; recall +0.0440 [0.0280, 0.0590] \\
\midrule
HotpotQA & 50k & Question only & 0.9820 & 0.7480 & 0.8650 & -- & -- \\
HotpotQA & 50k & Single rewritten query & 0.9720 & 0.6520 & 0.8120 & -- & -- \\
HotpotQA & 50k & Question + hypothetical passage & 0.9860 & \best{0.8560} & \best{0.9210} & All +0.1080 [0.0680, 0.1460]; recall +0.0560 [0.0350, 0.0770] & All +0.2040 [0.1620, 0.2440]; recall +0.1090 [0.0860, 0.1320] \\
\midrule
HotpotQA & 100k & Question only & 0.9800 & 0.6960 & 0.8380 & -- & -- \\
HotpotQA & 100k & Single rewritten query & 0.9680 & 0.5960 & 0.7820 & -- & -- \\
HotpotQA & 100k & Question + hypothetical passage & 0.9780 & \best{0.8340} & \best{0.9060} & All +0.1380 [0.1000, 0.1780]; recall +0.0680 [0.0470, 0.0890] & All +0.2380 [0.1940, 0.2800]; recall +0.1240 [0.0990, 0.1480] \\
\midrule
2Wiki & Local & Question only & 0.9980 & 0.5100 & 0.7655 & -- & -- \\
2Wiki & Local & Single rewritten query & 1.0000 & 0.4720 & 0.7380 & -- & -- \\
2Wiki & Local & Question + hypothetical passage & 1.0000 & \best{0.7800} & \best{0.9130} & All +0.2700 [0.2300, 0.3100]; recall +0.1475 [0.1290, 0.1675] & All +0.3080 [0.2680, 0.3520]; recall +0.1750 [0.1545, 0.1960] \\
\midrule
2Wiki & 50k & Question only & 1.0000 & 0.4460 & 0.7270 & -- & -- \\
2Wiki & 50k & Single rewritten query & 0.9960 & 0.4040 & 0.6900 & -- & -- \\
2Wiki & 50k & Question + hypothetical passage & 0.9980 & \best{0.6900} & \best{0.8655} & All +0.2440 [0.2040, 0.2840]; recall +0.1385 [0.1175, 0.1595] & All +0.2860 [0.2440, 0.3280]; recall +0.1755 [0.1545, 0.1980] \\
\midrule
2Wiki & 100k & Question only & 1.0000 & 0.4180 & 0.7105 & -- & -- \\
2Wiki & 100k & Single rewritten query & 0.9940 & 0.3600 & 0.6620 & -- & -- \\
2Wiki & 100k & Question + hypothetical passage & 0.9980 & \best{0.6500} & \best{0.8440} & All +0.2320 [0.1900, 0.2760]; recall +0.1335 [0.1110, 0.1560] & All +0.2900 [0.2460, 0.3340]; recall +0.1820 [0.1585, 0.2055] \\
\bottomrule
\end{tabularx}
\begin{tablenotes}
\footnotesize
\item Local denotes the original retained-duplicate local index used in the main stronger-encoder reproduction table (4,977 HotpotQA documents and 5,000 2Wiki documents). The rounded 5k stress-test cache is not used as a local-index replicate because the stress-construction script exact-deduplicates title--paragraph records before expansion and, for HotpotQA, pads the 4,977-document local corpus to 5,000 with train distractors. The 50k and 100k corpora preserve all released test-corpus records and add exact-title--paragraph-deduplicated contexts sampled from the corresponding IRCoT processed training split. All rows use \texttt{BAAI/bge-base-en-v1.5}, CLS pooling, 512-token dense inputs, top-$k=10$, and the frozen generated query artifacts used in the stronger-encoder reproduction. Deltas are computed only for the HyDE query row and use paired bootstrap 95\% confidence intervals over 500 per-example retrieval scores. The summary CSV is \texttt{local-artifacts/corpus\_scale\_stress/corpus\_scale\_retrieval\_summary.csv}; Local retrieval JSONL files are the stronger-encoder reproduction artifacts in \texttt{local-artifacts/results}, while expanded-corpus retrieval JSONL files and embedding caches are stored under \texttt{local-artifacts/corpus\_scale\_stress}. This gold-preserving distractor expansion is a diagnostic stress test and should not be interpreted as full-wiki retrieval.
\end{tablenotes}
\end{threeparttable}
\end{table*}


\clearpage
\section{Answer and Evidence Metric Normalization}
\label{sec:supp_metric_normalization}

The answer scores reported as EM and token-level F1 use the same local evaluator for HotpotQA and 2WikiMultihopQA. This table defines answer scoring; it is separate from the answer-overlap diagnostics in Section~\ref{sec:supp_answer_overlap_normalization}, which compare generated query text against gold answers rather than scoring reader predictions.

\begin{table*}[t]
\centering
\caption{Answer EM/F1 normalization and scoring rules.}
\label{tab:answer_metric_definition}
\scriptsize
\begin{threeparttable}
\setlength{\tabcolsep}{4pt}
\begin{tabularx}{\textwidth}{>{\raggedright\arraybackslash}p{0.25\textwidth}
                            >{\raggedright\arraybackslash}X}
\toprule
Item & Rule used for reported answer metrics \\
\midrule
Evaluator & Local HotpotQA-style reimplementation in \texttt{src/evaluation/evaluate\_answers.py}; SHA-256 prefix \texttt{55A65166A2FB}. The official HotpotQA script is not imported at runtime, but the normalization steps match the standard HotpotQA answer-normalization recipe. \\
Shared across datasets & The same evaluator functions, \texttt{exact\_match} and \texttt{f1\_score}, are used for HotpotQA and 2WikiMultihopQA rows. No dataset-specific answer normalizer is used. \\
Lowercasing & Convert prediction and gold answer to Python strings and lowercase them before comparison. \\
Article removal & Remove standalone English articles \texttt{a}, \texttt{an}, and \texttt{the} using the regular expression \texttt{\textbackslash b(a|an|the)\textbackslash b}. \\
Punctuation & Delete ASCII punctuation characters from Python \texttt{string.punctuation}. Unicode punctuation is not separately canonicalized. \\
Whitespace & Split on whitespace and rejoin with a single ASCII space; leading, trailing, and repeated whitespace therefore collapse after lowercasing, article removal, and punctuation deletion. \\
Exact match & EM is 1 if the normalized prediction string is exactly equal to the normalized gold string, and 0 otherwise. \\
Token F1 & Tokenize the normalized strings by whitespace. Let $c$ be the multiset intersection count between predicted and gold tokens. If either side is empty and the other is nonempty, F1 is 0; if both are empty, F1 is 1. Otherwise precision is $c/|\mathrm{pred}|$, recall is $c/|\mathrm{gold}|$, and F1 is $2PR/(P+R)$. \\
Multiple gold answers & The released answer rows store one \texttt{gold\_answer} string per example. The evaluator does not take a maximum over multiple aliases or alternative gold answers. \\
Empty predictions & Empty predictions are normalized to the empty string. They receive EM 1 and F1 1 only against an empty normalized gold answer; otherwise EM and F1 are 0. \\
Yes/no, numbers, dates & No special handling is applied beyond the same lowercasing, article removal, ASCII punctuation deletion, and whitespace collapse. Number words are not converted to digits, dates are not normalized, and yes/no answers are ordinary strings. \\
\bottomrule
\end{tabularx}
\begin{tablenotes}
\footnotesize
\item The same functions are imported by the answer-summary and paired-comparison scripts, including the BGE query-composition and equal-budget reader diagnostics.
\end{tablenotes}
\end{threeparttable}
\end{table*}

\begin{table*}[t]
\centering
\caption{Supporting-title and supporting-paragraph metric normalization rules.}
\label{tab:supporting_evidence_metric_definition}
\scriptsize
\begin{threeparttable}
\setlength{\tabcolsep}{4pt}
\begin{tabularx}{\textwidth}{>{\raggedright\arraybackslash}p{0.25\textwidth}
                            >{\raggedright\arraybackslash}X}
\toprule
Item & Rule used for reported evidence metrics \\
\midrule
Main supporting-title scoring & The reported any-hit@10, all-support hit@10, and supporting-title recall@10 use exact title strings from the retrieval artifact: \texttt{set(row["supporting\_facts"]["title"])} for gold titles and \texttt{\{doc["title"] for doc in row["retrieved"]\}} for retrieved titles. \\
Title lowercasing and trimming & Main title metrics do not lowercase titles, strip titles, or delete ASCII/Unicode punctuation at scoring time. They rely on the released processed-data title strings and retrieved document title strings matching exactly. \\
Duplicate titles & Gold and retrieved titles are converted to Python sets before scoring, so repeated retrieved records with the same title count once. A row is all-support hit@10 when the nonempty gold-title set is a subset of the retrieved-title set; recall is the intersection size divided by the gold-title set size. \\
Evaluator paths and hashes & The canonical title evaluator is \texttt{src/evaluation/evaluate\_retrieval.py}, SHA-256 prefix \texttt{CDBB6B75E260}. Equal-budget title contrasts use the same set rule in \texttt{src/evaluation/analyze\_equal\_budget\_pairwise.py}, SHA-256 prefix \texttt{46AD7C79BC18}; BGE retrieval-generation scripts store the same metrics during retrieval. \\
Paragraph-audit normalization & The supporting-paragraph audit uses \texttt{normalize\_text(value) = " ".join(str(value or "").split()).casefold()} for both titles and paragraph texts, then forms exact keys as \texttt{(normalized\_title, normalized\_paragraph\_text)}. It collapses whitespace and casefolds Unicode case, but it does not remove ASCII or Unicode punctuation. \\
Paragraph-audit gold keys & Gold paragraph keys are derived from corpus records with \texttt{is\_supporting=true} and the matching \texttt{source\_question\_id}. Retrieved paragraph keys are built from the top-10 retrieved documents using the same normalized title--paragraph-text function. \\
Paragraph-audit duplicates & Paragraph keys, record ids, and titles are stored as sets. Duplicate title--paragraph records therefore collapse for paragraph-key scoring, while the optional record-id metric remains stricter because it requires the source-question-specific \texttt{doc\_id}. \\
Paragraph-audit script & The audit script is \texttt{src/evaluation/analyze\_supporting\_paragraph\_metrics.py}, SHA-256 prefix \texttt{F8200E6D0911}. Its generated table is Table~\ref{tab:supporting_paragraph_audit}. \\
\bottomrule
\end{tabularx}
\begin{tablenotes}
\footnotesize
\item Title-level metrics are the primary retrieval metrics in the main manuscript because the benchmark annotations are supporting-title based in the released processed rows. Paragraph-level metrics are a stricter artifact audit over the local corpus records, not a new sentence-level supporting-fact metric.
\end{tablenotes}
\end{threeparttable}
\end{table*}

\clearpage
\section{Answer-Overlap Normalization}
\label{sec:supp_answer_overlap_normalization}

The following tables support the answer-string-overlap boundary and the secondary 2Wiki generalization analysis. The overlap diagnostics are query-side diagnostics: the hypothetical passage is used for retrieval, while the reader receives only the original question and retrieved corpus evidence.

\begin{table*}[t]
\centering
\caption{Reproducibility rules for answer-string-overlap diagnostics.}
\label{tab:answer_overlap_definition}
\scriptsize
\begin{threeparttable}
\setlength{\tabcolsep}{4pt}
\begin{tabularx}{\textwidth}{>{\raggedright\arraybackslash}p{0.24\textwidth}
                            >{\raggedright\arraybackslash}X}
\toprule
Item & Rule used in the released audit scripts \\
\midrule
Normalization & Convert to lowercase; remove standalone English articles \texttt{a}, \texttt{an}, and \texttt{the}; remove ASCII punctuation using Python \texttt{string.punctuation}; collapse consecutive whitespace. Punctuation is deleted rather than replaced, so hyphens, slashes, apostrophes, periods, and commas are removed under the same rule. \\
Overlap test & Let $g$ be the normalized gold answer and $z$ be the normalized generated query text: the HyDE hypothetical passage for HyDE rows or the rewritten query for the matched-call rewrite row. An overlap is counted when $g$ is nonempty, $z$ is nonempty, and $g$ occurs as an exact substring of $z$. \\
Alias handling & No alias dictionary, entity linking, date parser, stemming, or paraphrase matching is used. Aliases count only if the normalized benchmark gold string itself appears in the generated query text. \\
Numbers and dates & No special numeric or date canonicalization is applied beyond punctuation removal and whitespace normalization. Number words are not mapped to digits, dates are not reordered, hyphenated dates lose their hyphens, and partial date aliases are not counted unless they satisfy the exact normalized substring rule. \\
Stop-answer or short-answer rule & A gold answer is flagged as short/ambiguous if its normalized form is \texttt{yes}, \texttt{no}, a pure numeric string, or a single token with at most three characters. This is the only stop-answer/minimum-length filter used for the non-trivial overlap count. \\
Non-trivial answer overlap & A non-trivial overlap is an overlap where the gold answer is not flagged as short/ambiguous. Thus the HotpotQA count of 252 is the subset of the 292 overlap cases after removing short/ambiguous answers. \\
Answer-absent subset & The answer-absent subset contains examples with no exact normalized overlap between the gold answer and the generated HyDE hypothetical passage. For HotpotQA, this gives $500-292=208$ examples; for 2WikiMultihopQA, this gives $500-347=153$ examples. \\
\bottomrule
\end{tabularx}
\begin{tablenotes}
\footnotesize
\item The scripts implementing these rules are \texttt{analyze\_hyde\_leakage.py}, \texttt{analyze\_query\_answer\_overlap.py}, and \texttt{analyze\_answer\_absent\_subset.py}. The rule is intentionally lexical: it is reproducible, but it does not rule out paraphrased answer content, aliases, model memory, or benchmark contamination.
\end{tablenotes}
\end{threeparttable}
\end{table*}

The exact lexical implementation used by the released audit scripts is:
\begingroup\footnotesize
\begin{verbatim}
def norm(s):
    s = lower(str(s))
    s = remove_regex_word_articles(s, articles={"a", "an", "the"})
    s = delete_ascii_punctuation(s, chars=string.punctuation)
    return collapse_whitespace(s)

def is_short_or_ambiguous(gold):
    g = norm(gold)
    return (
        g in {"yes", "no"}
        or regex_fullmatch(r"\d+(?:\.\d+)?", g)
        or (len(split(g)) == 1 and len(g) <= 3)
    )

def has_answer_string_overlap(gold, generated_query):
    g = norm(gold)
    z = norm(generated_query)
    return bool(g) and bool(z) and g in z

def is_non_trivial_overlap(gold, generated_query):
    return has_answer_string_overlap(gold, generated_query) \
        and not is_short_or_ambiguous(gold)
\end{verbatim}
\endgroup
No alias table, entity linker, date parser, stemming, paraphrase matcher, or
numeric word-to-digit converter is applied before these functions.

\begin{table}[t]
\centering
\caption{HyDE hypothetical-passage answer-overlap diagnostic on HotpotQA.}
\label{tab:hyde_overlap_diagnostic}
\scriptsize
\begin{threeparttable}
\setlength{\tabcolsep}{3pt}
\begin{tabularx}{\columnwidth}{>{\raggedright\arraybackslash}X r r r}
\toprule
Group & $N$ & All hit@10 & Recall@10 \\
\midrule
All examples & 500 & 0.8680 & 0.9280 \\
Gold answer string appears & 292 & 0.9452 & 0.9726 \\
Gold answer string absent & 208 & 0.7596 & 0.8654 \\
Non-trivial answer string appears & 252 & 0.9365 & 0.9683 \\
\bottomrule
\end{tabularx}
\begin{tablenotes}
\footnotesize
\item A match is a normalized string diagnostic over generated hypothetical passages using the rules in Table~\ref{tab:answer_overlap_definition}. The hypothetical passage is used only as a retrieval query and is not included in the reader evidence prompt.
\end{tablenotes}
\end{threeparttable}
\end{table}

\begin{table*}[t]
\centering
\caption{Answer-overlap robustness analysis on HotpotQA.}
\label{tab:hyde_overlap_robustness}
\scriptsize
\begin{threeparttable}
\setlength{\tabcolsep}{3pt}
\begin{tabularx}{\textwidth}{>{\raggedright\arraybackslash}X r r r r r r c}
\toprule
Group & $N$ & Dense all-hit & HyDE all-hit & $\Delta$ all-hit & Dense F1 & HyDE F1 & W$\rightarrow$C / C$\rightarrow$W \\
\midrule
all examples & 500 & 0.6300 & 0.8680 & 0.2380 & 0.7199 & 0.8105 & 53 / 8 \\
answer in hyde & 292 & 0.6644 & 0.9452 & 0.2808 & 0.7976 & 0.9085 & 33 / 0 \\
answer not in hyde & 208 & 0.5817 & 0.7596 & 0.1779 & 0.6110 & 0.6730 & 20 / 8 \\
nontrivial answer in hyde & 252 & 0.6825 & 0.9365 & 0.2540 & 0.8130 & 0.8953 & 22 / 0 \\
\bottomrule
\end{tabularx}
\begin{tablenotes}
\footnotesize
\item Groups are defined by normalized string overlap between the gold answer and the generated hypothetical passage using the rules in Table~\ref{tab:answer_overlap_definition}. The answer-not-in-HyDE row is the robustness slice most relevant to answer-memory concerns; it uses existing answer files and does not require new LLM calls.
\end{tablenotes}
\end{threeparttable}
\end{table*}

\begin{table*}[t]
\centering
\caption{Answer-absent subset paired analysis.}
\label{tab:answer_absent_subset}
\scriptsize
\begin{threeparttable}
\setlength{\tabcolsep}{3pt}
\begin{tabularx}{\textwidth}{>{\raggedright\arraybackslash}X r r r r r r r c c c}
\toprule
Dataset & $N$ & Dense all-hit & HyDE all-hit & Dense recall & HyDE recall & Dense F1 & HyDE F1 & $\Delta$F1 [95\% CI] & W$\rightarrow$C / C$\rightarrow$W & McNemar $p$ \\
\midrule
HotpotQA & 208 & 0.5817 & 0.7596 & 0.7812 & 0.8654 & 0.6110 & 0.6730 & 0.0620 [0.0170, 0.1080] & 20 / 8 & 0.0357 \\
2WikiMultihopQA & 153 & 0.3856 & 0.6863 & 0.6797 & 0.8480 & 0.4058 & 0.6076 & 0.2019 [0.1372, 0.2671] & 30 / 3 & $<0.0001$ \\
\bottomrule
\end{tabularx}
\begin{tablenotes}
\footnotesize
\item The subset contains examples where the normalized gold-answer string is absent from the generated HyDE hypothetical passage under the rules in Table~\ref{tab:answer_overlap_definition}. All rows reuse frozen retrieval and reader answer artifacts. The diagnostic reduces the exact answer-string overlap explanation, but it does not rule out paraphrased answer content, benchmark contamination, or model memorization.
\end{tablenotes}
\end{threeparttable}
\end{table*}


The following gold-content diagnostics refine the exact-answer overlap analysis. They are not used for model selection or inference. The first table assigns each hypothetical passage to a mutually exclusive automatic stratum: exact answer string, alias/paraphrase proxy, supporting-title/entity mention, or no identifiable gold content under the implemented lexical/entity rules. The second table performs a gold-aware offline retrieval-only ablation by deleting exact answer strings, and then deleting exact answer strings plus gold supporting-title strings, from the generated hypothetical passage before retrieval. These scrubbed rows are diagnostic only because gold answers and support titles are unknown at test time.

\begin{table*}[t]
\centering
\caption{Automatic gold-content stratification of HyDE hypothetical passages.}
\label{tab:hyde_gold_content_groups}
\scriptsize
\begin{threeparttable}
\setlength{\tabcolsep}{1.8pt}
\begin{tabularx}{\textwidth}{>{\raggedright\arraybackslash}p{0.10\textwidth} >{\raggedright\arraybackslash}X r r r r r r r r r}
\toprule
Dataset & Stratum & $N$ & Dense all & HyDE all & $\Delta$ all [95\% CI] & Dense recall & HyDE recall & Dense F1 & HyDE F1 & $\Delta$F1 [95\% CI] \\
\midrule
HotpotQA & Exact-answer present & 292 & 0.6644 & 0.9452 & 0.2808 [0.2295, 0.3322] & 0.8271 & 0.9726 & 0.7976 & 0.9085 & 0.1110 [0.0756, 0.1461] \\
HotpotQA & Alias/paraphrase proxy present & 22 & 0.6818 & 0.9545 & 0.2727 [0.0909, 0.4545] & 0.8409 & 0.9773 & 0.7072 & 0.7841 & 0.0768 [-0.0379, 0.2078] \\
HotpotQA & Supporting-entity present & 154 & 0.6104 & 0.7987 & 0.1883 [0.1234, 0.2532] & 0.8019 & 0.8961 & 0.6238 & 0.6786 & 0.0548 [-0.0006, 0.1057] \\
HotpotQA & No identifiable gold content & 32 & 0.3750 & 0.4375 & 0.0625 [0.0000, 0.1562] & 0.6406 & 0.6406 & 0.4828 & 0.5693 & 0.0865 [-0.0312, 0.2167] \\
2Wiki & Exact-answer present & 347 & 0.3458 & 0.6599 & 0.3141 [0.2680, 0.3631] & 0.6650 & 0.8617 & 0.6520 & 0.7824 & 0.1304 [0.0916, 0.1738] \\
2Wiki & Alias/paraphrase proxy present & 20 & 0.5500 & 0.9500 & 0.4000 [0.1500, 0.6500] & 0.7750 & 0.9750 & 0.5375 & 0.7786 & 0.2411 [0.0571, 0.4500] \\
2Wiki & Supporting-entity present & 132 & 0.3636 & 0.6515 & 0.2879 [0.2121, 0.3712] & 0.6667 & 0.8314 & 0.3889 & 0.5863 & 0.1974 [0.1304, 0.2642] \\
2Wiki & No identifiable gold content & 1 & 0.0000 & 0.0000 & 0.0000 [0.0000, 0.0000] & 0.5000 & 0.5000 & 0.0000 & 0.0000 & 0.0000 [0.0000, 0.0000] \\
\bottomrule
\end{tabularx}
\begin{tablenotes}
\footnotesize
\item The strata are mutually exclusive and assigned in the order shown: exact normalized gold answer, automatic alias/paraphrase proxy, supporting-title/entity mention, and no identifiable gold content. The proxy is lexical and conservative; it does not replace human leakage annotation.
\end{tablenotes}
\end{threeparttable}
\end{table*}

\begin{table*}[t]
\centering
\caption{Gold-aware scrubbed HyDE retrieval-only diagnostic.}
\label{tab:hyde_gold_content_scrubbed}
\scriptsize
\begin{threeparttable}
\setlength{\tabcolsep}{4pt}
\begin{tabularx}{\textwidth}{>{\raggedright\arraybackslash}p{0.14\textwidth} >{\raggedright\arraybackslash}X r r r r}
\toprule
Dataset & Retrieval query & All hit@10 & Recall@10 & $\Delta$ all vs Dense [95\% CI] & $\Delta$ all vs original HyDE [95\% CI] \\
\midrule
HotpotQA & Dense question only & 0.6300 & 0.8080 & 0.0000 [0.0000, 0.0000] & -0.2380 [-0.2760, -0.2000] \\
HotpotQA & Original HyDE q+h & 0.8680 & 0.9280 & 0.2380 [0.2020, 0.2780] & 0.0000 [0.0000, 0.0000] \\
HotpotQA & Answer-scrubbed HyDE q+h & 0.8380 & 0.9120 & 0.2080 [0.1720, 0.2460] & -0.0300 [-0.0460, -0.0160] \\
HotpotQA & Answer+support-title-scrubbed HyDE q+h & 0.7800 & 0.8830 & 0.1500 [0.1140, 0.1860] & -0.0880 [-0.1140, -0.0620] \\
2Wiki & Dense question only & 0.3580 & 0.6695 & 0.0000 [0.0000, 0.0000] & -0.3100 [-0.3520, -0.2700] \\
2Wiki & Original HyDE q+h & 0.6680 & 0.8575 & 0.3100 [0.2700, 0.3520] & 0.0000 [0.0000, 0.0000] \\
2Wiki & Answer-scrubbed HyDE q+h & 0.6480 & 0.8445 & 0.2900 [0.2500, 0.3340] & -0.0200 [-0.0340, -0.0060] \\
2Wiki & Answer+support-title-scrubbed HyDE q+h & 0.5640 & 0.7770 & 0.2060 [0.1660, 0.2480] & -0.1040 [-0.1320, -0.0780] \\
\bottomrule
\end{tabularx}
\begin{tablenotes}
\footnotesize
\item The scrubbed rows are offline diagnostics that delete gold-aware strings from the generated hypothetical passage before retrieval. They are not deployable inference methods because gold answers and supporting titles are unavailable at test time. Reader answers are not regenerated for scrubbed rows.
\end{tablenotes}
\end{threeparttable}
\end{table*}

\begin{table*}[t]
\centering
\caption{Retrieval-only HyDE query-composition check on the released IRCoT 2WikiMultihopQA \texttt{test\_subsampled} split.}
\label{tab:2wiki_hyde_query_mechanism}
\small
\begin{threeparttable}
\begin{tabularx}{\textwidth}{>{\raggedright\arraybackslash}X r r r r}
\toprule
Query used for dense retrieval & $N$ & Any hit@10 & All hit@10 & Recall@10 \\
\midrule
Question only & 500 & 0.9840 & 0.3580 & 0.6695 \\
Hypothetical passage only & 500 & \best{0.9940} & 0.6640 & 0.8530 \\
Question + hypothetical passage & 500 & \best{0.9940} & \best{0.6680} & \best{0.8575} \\
\bottomrule
\end{tabularx}
\begin{tablenotes}
\footnotesize
\item This no-new-LLM diagnostic reuses the existing 2Wiki HyDE generations and dense index, without additional reader calls. Complete supporting-title retrieval is higher for the hypothetical-only row than for question-only retrieval, while keeping the original question provides a small additional anchor.
\end{tablenotes}
\end{threeparttable}
\end{table*}


\clearpage
\section{Query-Length Matched HyDE Sensitivity}
\label{sec:supp_query_length_matched_hyde}

The matched-call query-reformulation control equalizes query-side LLM invocation count but not generated-text length or information budget. The following no-new-LLM audit reports MiniLM tokenizer lengths for the main query-composition variants. For the hypothetical-only control, the HyDE hypothetical passage is truncated to the corresponding rewritten-query token count. For the question-plus-hypothetical control, only the hypothetical passage is truncated, and the serialized query is matched to the corresponding question-plus-rewrite token length. The paired deltas use the comparator shown in Table~\ref{tab:query_length_matched_hyde}: the hypothetical-only row is compared with the single rewritten query, whereas the question-plus-hypothetical row is compared with the question-plus-rewritten-query control. This analysis is a query-composition sensitivity check: it does not replace the main HyDE result, but it shows that matched-length rewriting alone does not recover the full HyDE result.

\begin{table*}[t]
\centering
\caption{Query-token length audit and length-matched HyDE retrieval sensitivity on HotpotQA. Token counts use the MiniLM tokenizer. The hypothetical-only control matches the rewritten-query length; the question-plus-hypothetical control truncates only the hypothetical passage so that the serialized query is length-matched to question plus rewrite. No new LLM calls are made.}
\label{tab:query_length_matched_hyde}
\scriptsize
\begin{threeparttable}
\setlength{\tabcolsep}{3pt}
\begin{tabularx}{\textwidth}{>{\raggedright\arraybackslash}Xrrrr}
\toprule
Query & Mean tokens & Median & P95 & Hits 256-token cap \\
\midrule
q & 19.9140 & 19.0000 & 33 & 0 \\
r & 10.9380 & 10.0000 & 17 & 0 \\
q+r & 35.8520 & 34.0000 & 52 & 0 \\
h & 68.7400 & 68.0000 & 92 & 0 \\
q+h & 92.6540 & 92.0000 & 119 & 0 \\
\bottomrule
\end{tabularx}
\vspace{1mm}
\begin{tabularx}{\textwidth}{>{\raggedright\arraybackslash}X>{\raggedright\arraybackslash}Xrrrrr}
\toprule
Comparator & Target retrieval query & Any-hit@10 & All-hit@10 & Recall@10 & $\Delta$ All-hit@10 & $\Delta$ Recall@10 \\
\midrule
Single rewritten query & Length-matched hypothetical only & 0.9200 & 0.2740 & 0.5970 & -0.3440 [-0.3940, -0.2980] & -0.2020 [-0.2310, -0.1730] \\
Question + rewritten query & Length-matched question + hypothetical & 0.9740 & 0.6400 & 0.8070 & 0.0120 [-0.0240, 0.0480] & 0.0040 [-0.0170, 0.0250] \\
\bottomrule
\end{tabularx}
\begin{tablenotes}
\footnotesize
\item Deltas are paired against the comparator shown in the first column: the hypothetical-only control is compared with the single rewritten query, and the question-plus-hypothetical control is compared with the question-plus-rewritten-query control. The question-plus-hypothetical row matches the serialized query length to the corresponding question-plus-rewrite query, up to tokenizer discreteness, while preserving the original question as an anchor in this length-matched setting. The sensitivity controls show that matched-length rewriting does not recover the full HyDE result, so the safest interpretation is a combined effect of document-like generation and a larger semantic information budget rather than document-like form alone.
\end{tablenotes}
\end{threeparttable}
\end{table*}


\clearpage
\section{Error Taxonomy Rules}
\label{sec:supp_error_taxonomy_rules}

The annotation unit for the error taxonomy is one exact-match mismatch from the final HyDE-style RAG plus guarded verifier HotpotQA run. Each mismatch is assigned to one primary diagnostic category using frozen retrieval and answer artifacts: retrieval failure if no required supporting title is retrieved; partial evidence or missing hop if at least one but not all supporting titles are retrieved; reader reasoning failure if all supporting titles are retrieved but the annotator judged the answer semantically wrong; and surface-form, alias, or granularity mismatch if the annotator judged the retrieved evidence to support an acceptable answer that benchmark normalization still rejects by surface form or granularity.

When multiple explanations are plausible, the precedence order is retrieval availability, then reader reasoning, then format or alias. Ambiguous cases are resolved conservatively as not verifier-fixable. The annotations are used only for interpretation and case selection, not for training, prompting, or model selection. The manual inspection stage was performed by a single annotator, with no independent second-pass adjudication or inter-annotator agreement estimate, so the taxonomy should be treated as a descriptive diagnostic rather than a formally validated annotation benchmark.

\begin{table}[t]
\centering
\caption{Single-annotator EM-mismatch diagnostic on HotpotQA.}
\label{tab:single_annotator_error_taxonomy}
\small
\begin{threeparttable}
\begin{tabularx}{\textwidth}{>{\raggedright\arraybackslash}X r r}
\toprule
Diagnostic category & Count & Share of 156 EM mismatches \\
\midrule
Retrieval failure & 6 & 3.8\% \\
Partial evidence / missing hop & 39 & 25.0\% \\
Reader reasoning failure despite complete evidence & 20 & 12.8\% \\
Annotator-judged surface-form, alias, or granularity mismatch & 91 & 58.3\% \\
\bottomrule
\end{tabularx}
\begin{tablenotes}
\footnotesize
\item These counts describe residual exact-match mismatches in the final HyDE-style RAG + guarded verifier run. They come from a single-annotator diagnostic and are not a validated human-annotation benchmark.
\end{tablenotes}
\end{threeparttable}
\end{table}

\clearpage
\section{Corpus, Supporting-Title Completeness, and Case Artifacts}
\label{sec:supp_evidence_case_artifacts}

The pooled supporting-title completeness rows are descriptive method-example observations and are not treated as statistically independent samples. The qualitative cases were selected from rule-filtered candidates using frozen per-example artifacts and are intended to make the retrieval-centered evidence changes and downstream answer-format boundaries inspectable. Table~\ref{tab:qualitative_cases} gives the compact reader-facing summary, and Table~\ref{tab:qualitative_case_artifacts} exposes the corresponding generated passage, retrieved titles, evidence excerpts, reader answer, verifier boundary when invoked, and guarded final answer. The supporting-paragraph audit at the end of this section rescores the same retrieval files with exact normalized title--paragraph-text keys derived from \texttt{is\_supporting=true} corpus records.

\begin{table*}[t]
\centering
\caption{Local candidate-corpus statistics for the released processed splits.}
\label{tab:corpus_statistics}
\small
\begin{threeparttable}
\setlength{\tabcolsep}{5pt}
\begin{tabularx}{\textwidth}{>{\raggedright\arraybackslash}X r r}
\toprule
Statistic & HotpotQA & 2WikiMultihopQA \\
\midrule
Questions & 500 & 500 \\
Total candidate paragraphs & 4977 & 5000 \\
Unique titles & 4931 & 3273 \\
Mean / median source-pool records contributed per question & 10.0 / 10 & 10 / 10 \\
Source-pool records contributed per question, min--max & 3--10 & 10--10 \\
Mean gold supporting titles & 2.00 & 2.42 \\
Mean supporting paragraphs & 2.00 & 2.42 \\
Mean / median distractor paragraphs & 8.0 / 8 & 7.6 / 8 \\
Corpus organization & Joint index over local pools & Joint index over local pools \\
\bottomrule
\end{tabularx}
\begin{tablenotes}
\footnotesize
\item Source-pool records are the records contributed by each released question-level candidate pool before all records are unioned into a dataset-level joint index. Retrieval is performed over the full joint corpus, not over a separate per-question pool. Exact duplicate title--paragraph records are retained before indexing; if the same paragraph appears in multiple question-level pools, each occurrence remains a separate document with a source-question-specific document identifier. Dense, BM25, and hybrid retrieval encode only document titles and paragraph texts. Question identifiers, supporting-evidence flags, gold answers, and evaluation metadata are retained in JSONL artifacts for analysis but are excluded from retrieval representations.
\item In both released processed splits, each annotated supporting title is associated with exactly one gold supporting paragraph record for every question, which explains why the mean supporting-title and supporting-paragraph counts are identical.
\item Supporting-title metrics collapse duplicate retrieved titles before computing any-hit@10, all-hit@10, and supporting-title recall@10.
\end{tablenotes}
\end{threeparttable}
\end{table*}

\begin{table*}[t]
\centering
\caption{Supporting-title completeness and reader accuracy on the released IRCoT HotpotQA \texttt{test\_subsampled} split.}
\label{tab:evidence_completeness_buckets_supp}
\small
\begin{threeparttable}
\begin{tabularx}{\textwidth}{>{\raggedright\arraybackslash}X l r r r r}
\toprule
Scope & Supporting-title bucket & $N$ & EM & F1 & Mean title recall \\
\midrule
All retrieval-only method-example observations & full title support & 2181 & \best{0.7240} & \best{0.8584} & 1.0000 \\
All retrieval-only method-example observations & partial title support & 780 & 0.3846 & 0.4744 & 0.5000 \\
All retrieval-only method-example observations & no title support & 39 & 0.0769 & 0.0872 & 0.0000 \\
\midrule
HyDE-style RAG only & full title support & 434 & \best{0.7396} & \best{0.8720} & 1.0000 \\
HyDE-style RAG only & partial title support & 60 & 0.3500 & 0.4472 & 0.5000 \\
HyDE-style RAG only & no title support & 6 & 0.0000 & 0.0000 & 0.0000 \\
\bottomrule
\end{tabularx}
\begin{tablenotes}
\footnotesize
\item Full title support means all annotated supporting titles are retrieved in the top-10 evidence set; partial title support means at least one but not all supporting titles are retrieved. The pooled rows are descriptive method-example observations across the six retrieval-only methods and therefore are not treated as independent examples.
\end{tablenotes}
\end{threeparttable}
\end{table*}

\begin{table*}[t]
\centering
\caption{Reproducibility checklist and canonical artifact ledger.}
\label{tab:reproducibility_checklist}
\scriptsize
\begin{threeparttable}
\setlength{\tabcolsep}{3pt}
\begin{tabularx}{\textwidth}{>{\raggedright\arraybackslash}p{0.16\textwidth}
                            >{\raggedright\arraybackslash}p{0.31\textwidth}
                            >{\raggedright\arraybackslash}X}
\toprule
Component & Setting reported in the paper & Audit artifact or reproducibility check \\
\midrule
Data splits & Released IRCoT processed HotpotQA and 2WikiMultihopQA \texttt{test\_subsampled} splits; 500 examples per split. & Upstream repository, commit, source filenames, local download timestamps, raw SHA-256 hashes, converted question/corpus JSONL SHA-256 hashes, and conversion-script hash are listed in Table~\ref{tab:dataset_artifact_provenance}. The processed local-corpus protocol is documented in \texttt{docs/reproducibility\_protocol.md}. \\
Evidence pool & Local candidate evidence pools from the released processed data; paragraph-level documents with title plus paragraph text. & Corpus statistics are reported in the main manuscript and generated by \texttt{src/evaluation/analyze\_corpus\_statistics.py}; results are not reported as full-wiki IRCoT reproduction. \\
Retrieval settings & Top-$k=10$ final evidence documents; main dense retrieval uses \texttt{all-MiniLM-L6-v2}, normalized embeddings, and FAISS inner-product search; post-hoc BGE-base diagnostics use the shared configuration in Section~\ref{sec:supp_dense_encoder_truncation}. & Retriever scripts in \texttt{src/retriever}; dense, BM25, hybrid, multi-query, iterative, HyDE, BGE query-composition, and equal-budget BGE retrieval outputs are retained in \texttt{local-artifacts/results} and \texttt{local-artifacts}. Multi-query and iterative rows are heuristic controls, not faithful public-method reproductions. \\
Prompt settings & HyDE generation, evidence-grounded short-answer reading, and conservative verification use fixed prompt templates; reader/verifier evidence blocks are capped at 900 characters per document. & Prompt builders in \texttt{src/generator} and \texttt{src/verifier}; prompt JSONL rows store metadata for evaluation, but gold answers are not inserted into the model prompt text. \\
LLM settings & Hosted Qwen3.7-Max through the Alibaba Cloud Bailian/DashScope OpenAI-compatible chat-completions interface with \texttt{temperature=0.0}. Reported primary reader, query-generation, and verifier artifacts use the API model string \texttt{qwen3.7-max}. The targeted equal-budget reader robustness check uses \texttt{qwen-turbo} over fixed retrieved evidence. The configured \texttt{max\_tokens} budgets are 64 for HyDE passage generation, rewritten-query generation, reader answers, and verifier JSON output; temperature 0 is not claimed to ensure exact replay. & Generation script records \texttt{model}, prompt, retrieved evidence, prediction, and gold-answer metadata in answer JSONL files; \texttt{--resume} prevents duplicate regeneration. The local invocation ledger is frozen in \texttt{docs/model\_invocation\_ledger.md}, and Table~\ref{tab:qwen_api_reproducibility} reports provider/interface, hosted-alias, decoding, retry, prompt-hash, and raw-output boundaries. \\
Canonical ledger & Main HotpotQA table uses the audited \texttt{qwen3.7-max} artifact ledger for reported reader comparisons. Rule-based Evidence-Expanded Iterative Retrieval is the fixed heuristic control reported in the main table; Single-query Reformulation is the matched-call mechanism control. & Canonical ledger: \texttt{local-artifacts/ircot\_test500\_qwen\_results.csv}; audit note: \texttt{local-artifacts/stage1\_experiment\_result\_audit\_2026-07-06.md}. Table~\ref{tab:configuration_provenance} separates the primary frozen study, post-hoc robustness diagnostics, and same-split verifier tuning. \\
Additional analyses & Corpus statistics, HyDE query-composition diagnostic, strict equal-budget query-form diagnostic, targeted Qwen-Turbo reader check, supporting-paragraph audit, top-$k$ sensitivity and worst-support-rank audit, corpus-scale retrieval stress test, corpus-scale leakage/hard-negative audit, answer-string-overlap diagnostic, paired reliability checks, evidence buckets, and efficiency accounting. & Stage notes and CSVs in \texttt{local-artifacts}; equal-budget artifacts are under \texttt{local-artifacts/equal\_budget\_query\_diagnostic}; the Qwen-Turbo paired reader statistics are summarized in \texttt{equal\_budget\_cross\_reader\_contrasts.csv} and \texttt{equal\_budget\_cross\_reader\_paired\_stats.md} by \texttt{src/evaluation/summarize\_equal\_budget\_reader.py}; the supporting-paragraph audit is under \texttt{local-artifacts/supporting\_paragraph\_audit} and generated by \texttt{src/evaluation/analyze\_supporting\_paragraph\_metrics.py}; top-$k$ artifacts are under \texttt{local-artifacts/topk\_sensitivity} and summarized by \texttt{src/evaluation/analyze\_topk\_sensitivity.py}; corpus-scale stress-test artifacts are under \texttt{local-artifacts/corpus\_scale\_stress} and generated by \texttt{src/evaluation/analyze\_corpus\_scale\_retrieval\_stress.py}; leakage-filtered and BGE-hard audit artifacts are under \texttt{local-artifacts/corpus\_scale\_audit} and generated by \texttt{src/evaluation/analyze\_corpus\_scale\_audit.py}. Corpus and efficiency tables are generated from actual JSONL artifacts rather than hand-estimated counts. \\
\bottomrule
\end{tabularx}
\begin{tablenotes}
\footnotesize
\item The checklist is intended to make the processed-data experimental protocol auditable. Artifact freezing records the exact files used for reporting, but it is not claimed as an untouched held-out development/test split protocol for every design choice. File paths are relative to the project root when possible. API credentials are loaded from local environment variables and are not part of the reproducibility artifact.
\end{tablenotes}
\end{threeparttable}
\end{table*}

\begin{table*}[t]
\centering
\caption{Artifact-grounded qualitative case studies emphasizing HyDE retrieval.}
\label{tab:qualitative_cases}
\scriptsize
\begin{threeparttable}
\setlength{\tabcolsep}{3pt}
\begin{tabularx}{\textwidth}{>{\raggedright\arraybackslash}p{0.17\textwidth}
                            >{\raggedright\arraybackslash}p{0.27\textwidth}
                            >{\raggedright\arraybackslash}p{0.23\textwidth}
                            >{\raggedright\arraybackslash}X}
\toprule
Case type & Question summary & Prediction / evidence state & Interpretation \\
\midrule
HyDE evidence acquisition & "Ew!" is a song by a television host born where? & Dense: United States; HyDE: Bay Ridge, Brooklyn; Gold: Bay Ridge, Brooklyn; Dense support partial (1/2); HyDE support full (2/2) & HyDE retrieves the missing supporting title without an exact normalized answer-string match in the hypothetical passage, turning a partial-evidence Dense failure into a correct answer. \\
HyDE evidence acquisition & How many films were directed by the director of Wise Blood? & Dense: I don't know; HyDE: 37; Gold: 37; Dense support partial (1/2), missing John Huston; HyDE support full (2/2) & HyDE retrieves the director page needed for the second hop and turns a partial-support Dense abstention into a correct count answer. \\
Answer-standardization boundary & Which star of Zork was also the voice of Pac-Man? & HyDE: Marty Ingels; Gold: Martin Ingerman; HyDE support full (2/2) & The reader returns the professional name Marty Ingels, whereas the benchmark gold answer uses the birth name Martin Ingerman. The two names refer to the same person, illustrating that exact-match evaluation can count semantically valid aliases as errors. \\
\bottomrule
\end{tabularx}
\begin{tablenotes}
\footnotesize
\item Cases were selected from rule-filtered candidates using frozen HotpotQA per-example artifacts: Dense and HyDE reader outputs, supporting-title states, and HyDE answer-overlap audit rows. They are illustrative rather than additional evaluation examples.
\end{tablenotes}
\end{threeparttable}
\end{table*}

\begin{table*}[t]
\centering
\caption{Compact artifact blocks for the qualitative case studies.}
\label{tab:qualitative_case_artifacts}
\scriptsize
\begin{threeparttable}
\setlength{\tabcolsep}{2pt}
\begin{tabularx}{\textwidth}{>{\raggedright\arraybackslash}p{0.15\textwidth}
                            >{\raggedright\arraybackslash}p{0.28\textwidth}
                            >{\raggedright\arraybackslash}p{0.28\textwidth}
                            >{\raggedright\arraybackslash}X}
\toprule
Example & Generated / retrieved artifacts & Evidence excerpts & Answer artifacts \\
\midrule
HyDE evidence acquisition\newline{\tiny\texttt{5ae68fcb5542992ae0d1635b}} & HyDE passage: "Ew!" is a 2014 comedy song performed by American television host Jimmy Fallon, featuring will.i.am. The track is based on a recurring sketch from The Tonight Show Starring Jimmy Fallon. Jimmy Fallon, the television host and comedian who created and perform...\newline Dense top: Ew!; Richie Kavanagh; Ezequiel Echeverría; supporting-title hits: Ew!; missing: Jimmy Fallon\newline HyDE top: Ew!; Jimmy Fallon; Tom Ewell; supporting-title hits: Ew!; Jimmy Fallon; missing: none & [Ew!] "Ew!" is a song by American television host and comedian Jimmy Fallon, featuring American rapper will.i.am. The song is based on a sketch from "The Tonight Show Starri...; [Jimmy Fallon] James Thomas Fallon (born September 19, 1974) is an American comedian, actor, television host, and musician. He is known for his work in television as a cast member on... & Initial reader answer: Bay Ridge, Brooklyn\newline Verifier boundary: not invoked for this example\newline Final guarded answer: Bay Ridge, Brooklyn \\
HyDE evidence acquisition\newline{\tiny\texttt{5ab3f6b9554299753aec5a03}} & HyDE passage: The 1979 film Wise Blood was directed by acclaimed American filmmaker John Huston. Over the course of his extensive career in Hollywood, Huston directed a total of 37 feature films, including celebrated classics such as The Maltese Falcon, The Treasure of t...\newline Dense top: Wise Blood (film); Wise Blood; Benedict Fitzgerald; supporting-title hits: Wise Blood (film); missing: John Huston\newline HyDE top: Wise Blood (film); John Huston; Benedict Fitzgerald; supporting-title hits: John Huston; Wise Blood (film); missing: none & [Wise Blood (film)] Wise Blood is an American 1979 drama film directed by John Huston and based on the 1952 novel "Wise Blood" by Flannery O'Connor. It was filmed mostly in and around Mac...; [John Huston] John Marcellus Huston ( ; August 5, 1906 – August 28, 1987) was an American film director, screenwriter and actor. He wrote the screenplays for most of the 37 feature... & Dense answer: I don't know; HyDE answer: 37.\newline Dense has partial support; HyDE has full support.\newline Verifier boundary: raw verifier proposed "37 feature films"; guard kept 37. \\
Answer-standardization boundary\newline{\tiny\texttt{5a728f015542991f9a20c4e4}} & HyDE passage: Marty Ingels, an American actor and voice artist best known for providing the original voice of the titular character in the 1982 animated television series Pac-Man, also starred in the live-action interactive video game Return to Zork.\newline Dense top: Marty Ingels; Mario Castañeda; David Acord; supporting-title hits: Marty Ingels; missing: Zork: Grand Inquisitor\newline HyDE top: Marty Ingels; Miłogost Reczek; Mario Castañeda; supporting-title hits: Marty Ingels; Zork: Grand Inquisitor; missing: none & [Marty Ingels] Martin Ingerman (March 9, 1936 – October 21, 2015), known professionally as Marty Ingels, was an American actor, comedian, comedy sketch writer and theatrical agent, w...; [Zork: Grand Inquisitor] Zork: Grand Inquisitor is a graphic adventure game developed by Activision and released for Windows in 1997, and for Macintosh in 2001. It builds upon the "Zork" and "... & Initial reader answer: Marty Ingels\newline Verifier boundary: not invoked for this example\newline Final guarded answer: Marty Ingels \\
\bottomrule
\end{tabularx}
\begin{tablenotes}
\footnotesize
\item Each block is a compact projection of frozen per-example artifacts. ``Not invoked'' means the verifier was not run for that example under the risk-selected verification protocol.
\end{tablenotes}
\end{threeparttable}
\end{table*}

\begin{table*}[t]
\centering
\caption{Supporting-paragraph audit for reported title-level retrieval rows.}
\label{tab:supporting_paragraph_audit}
\scriptsize
\begin{threeparttable}
\setlength{\tabcolsep}{3.5pt}
\begin{tabularx}{\textwidth}{>{\raggedright\arraybackslash}X l l r r r r}
\toprule
Dataset & Family & Method & Title all@10 & Paragraph all@10 & Paragraph recall@10 & Gap \\
\midrule
HotpotQA & Main & Dense & 0.6300 & 0.6300 & 0.8080 & +0.0000 \\
HotpotQA & Main & Multi-query & 0.6300 & 0.6300 & 0.8080 & +0.0000 \\
HotpotQA & Main & BM25 & 0.7320 & 0.7320 & 0.8590 & +0.0000 \\
HotpotQA & Main & Hybrid & 0.7480 & 0.7480 & 0.8690 & +0.0000 \\
HotpotQA & Main & Rule-based iterative & 0.7540 & 0.7540 & 0.8700 & +0.0000 \\
HotpotQA & Main & Direct rewrite & 0.6180 & 0.6180 & 0.7990 & +0.0000 \\
HotpotQA & Main & Question + rewrite & 0.6280 & 0.6280 & 0.8030 & +0.0000 \\
HotpotQA & Main & Hypothetical-only & 0.8480 & 0.8480 & 0.9180 & +0.0000 \\
HotpotQA & Main & HyDE & 0.8680 & 0.8680 & 0.9280 & +0.0000 \\
HotpotQA & Equal-budget & Keyword/entity & 0.9240 & 0.9240 & 0.9570 & +0.0000 \\
HotpotQA & Equal-budget & Direct rewrite & 0.8800 & 0.8800 & 0.9330 & +0.0000 \\
HotpotQA & Equal-budget & Decomposition & 0.8740 & 0.8740 & 0.9320 & +0.0000 \\
HotpotQA & Equal-budget & Document-like & 0.9260 & 0.9260 & 0.9600 & +0.0000 \\
2WikiMultihopQA & Main & Dense & 0.3580 & 0.3580 & 0.6695 & +0.0000 \\
2WikiMultihopQA & Main & BM25 & 0.4320 & 0.4320 & 0.7250 & +0.0000 \\
2WikiMultihopQA & Main & Hybrid & 0.4480 & 0.4480 & 0.7330 & +0.0000 \\
2WikiMultihopQA & Main & Hypothetical-only & 0.6640 & 0.6640 & 0.8530 & +0.0000 \\
2WikiMultihopQA & Main & HyDE & 0.6680 & 0.6680 & 0.8575 & +0.0000 \\
2WikiMultihopQA & Equal-budget & Keyword/entity & 0.6760 & 0.6760 & 0.8410 & +0.0000 \\
2WikiMultihopQA & Equal-budget & Direct rewrite & 0.4760 & 0.4760 & 0.7190 & +0.0000 \\
2WikiMultihopQA & Equal-budget & Decomposition & 0.4540 & 0.4540 & 0.7295 & +0.0000 \\
2WikiMultihopQA & Equal-budget & Document-like & 0.7520 & 0.7520 & 0.9035 & +0.0000 \\
\bottomrule
\end{tabularx}
\begin{tablenotes}
\footnotesize
\item The audited rows are exactly the rows listed in the table. They include the HotpotQA MiniLM main-result, query-composition, heuristic-control retrieval rows, the 2WikiMultihopQA reported MiniLM retrieval rows with available artifacts, and the four equal-budget query-form rows for both datasets. Title all@10 matches the main all-support hit@10 metric: all annotated supporting titles must appear in the top-10 retrieved evidence. Paragraph all@10 is stricter: all gold corpus records marked \texttt{is\_supporting=true} for the question must be recovered as exact normalized title--paragraph text. The gap is title all@10 minus paragraph all@10.
\item In the released processed splits, each annotated supporting title is associated with one supporting paragraph record, so identical title-level and paragraph-level values are expected when the retrieved title corresponds to the same processed paragraph text. The processed artifacts do not expose gold supporting sentence text in the 500-example JSONL questions used here, so this audit cannot compute sentence-level supporting-fact coverage. It primarily verifies exact processed-record recovery rather than introducing an independent sentence-level evidence notion.
\end{tablenotes}
\end{threeparttable}
\end{table*}


\clearpage
\section{Code and Artifact Availability}
\label{sec:supp_code_artifact_availability}

The anonymized review repository is available at:
\begin{center}
\href{https://github.com/creator-perterson/hyde-multihop-rag-reproducibility}{\texttt{https://github.com/creator-perterson/}\\\texttt{hyde-multihop-rag-reproducibility}}
\end{center}
It contains the implementation, prompt templates, frozen query and answer summaries, retrieval summaries, paired-analysis scripts, configuration-provenance records, model-invocation ledgers, answer-overlap diagnostics, and deduplication diagnostics, subject to the final venue policy and the licenses or terms of the upstream datasets and hosted-model provider. The release is intended to support artifact-level auditing and rescoring of the reported tables, rather than byte-identical replay of a hosted model alias.

The releasable artifact package will include local prompt text, retrieved-evidence serialization, local prediction text, parsed verifier JSON fields, stored model-string fields, artifact freeze timestamps, hashes, table-generation scripts, and aggregate analysis outputs. API credentials, private endpoint URLs, provider-side request identifiers, billing exports, latency logs, and provider raw request/response envelopes are not included because they are either private, were not retained by the local wrapper, or are not required for table-level rescoring. If venue, dataset, or hosted-provider terms restrict redistribution of generated raw text or third-party candidate-pool passages, the repository will instead provide file manifests, SHA-256 hashes, derived per-example statistics, parsed predictions/final answers where permitted, and scripts that reproduce the reported tables from authorized local copies.

\end{document}
